\chapter{SRAM、寄存器文件与 Cache}

\begin{keyidea}
离计算核心越近，存储越依赖低延迟、并行端口和局部性。本章从 SRAM 单元与阵列出发，建立寄存器文件和 Cache 的结构基础。
\end{keyidea}

\section{一、SRAM 单元与阵列}

SRAM（Static Random Access Memory）是 CPU 中速度最快、面积代价最高的存储器件。一个 SRAM bit 由 6 个 MOSFET 构成的双稳态锁存器保存，只要电源不断，数据就永远保持——不需要刷新，不需要时钟。代价是：每个 bit 需要 6 个晶体管，面积是 DRAM 的 6--8 倍。本节从单个 6T 单元的晶体管级结构讲起，经过读/写/保持操作的定量分析，到阵列组织、辅助电路、时序参数、功耗和软错误，最终给出真实芯片中的 SRAM 参数。

\begin{keyidea}
SRAM 的核心是\hwkey{正反馈锁存}：两个交叉耦合反相器形成双稳态，任何一端的微小扰动都会被反馈放大回稳态。这个正反馈既是 SRAM 保持数据的原因，也是写入时必须克服的障碍。所有设计权衡——读稳定性 vs 写能力、SNM vs 面积、assist 电路 vs 功耗——都源于这一对矛盾。
\end{keyidea}

%======================================================================
\topic{6T 单元结构：六个晶体管的角色}
%======================================================================

第〇节给出了从单点到系统的递进框架，本小节聚焦第一级：SRAM 的一个 bit 在晶体管级长什么样。\hwkey{6T SRAM 单元}由以下六个 MOSFET 组成：

\begin{longtable}{L{0.10\linewidth}L{0.12\linewidth}L{0.14\linewidth}L{0.54\linewidth}}
\toprule
编号 & 类型 & 角色 & 连接关系 \\
\midrule
P1 & PMOS & 上拉管（左反相器） & 源接 VDD，漏接 Q，栅接 $\overline{Q}$ \\\tablerowrule
N1 & NMOS & 下拉管（左反相器） & 源接 VSS，漏接 Q，栅接 $\overline{Q}$ \\\tablerowrule
P2 & PMOS & 上拉管（右反相器） & 源接 VDD，漏接 $\overline{Q}$，栅接 Q \\\tablerowrule
N2 & NMOS & 下拉管（右反相器） & 源接 VSS，漏接 $\overline{Q}$，栅接 Q \\\tablerowrule
M5 & NMOS & 访问管（左） & 源/漏接 BL 和 Q，栅接 WL \\\tablerowrule
M6 & NMOS & 访问管（右） & 源/漏接 $\overline{BL}$ 和 $\overline{Q}$，栅接 WL \\
\bottomrule
\end{longtable}

P1/N1 构成左反相器（输入 $\overline{Q}$，输出 Q），P2/N2 构成右反相器（输入 Q，输出 $\overline{Q}$）。两者交叉耦合形成\hwkey{双稳态锁存器}：Q=1/$\overline{Q}$=0 和 Q=0/$\overline{Q}$=1 是两个稳定状态。M5/M6 是\term{访问管}{access transistor}——控制锁存器与位线之间通断的开关管，栅极共同连接字线 WL；WL=0 时截止，锁存器与位线隔离，数据自保持；WL=1 时导通，位线与存储节点连通，既可以被读出（sense，把 Q/$\overline{Q}$ 的状态传到位线上），也可以被写驱动器反向驱动（把位线上的新数据强行写入存储节点）。

\begin{pdfdiagram}[x=1cm,y=1cm]
  % VDD rail
  \draw[BookTeal, line width=0.7pt] (1.0,3.2) -- (7.0,3.2);
  \node[BookTeal, font=\small\bfseries] at (4.0,3.5) {VDD};
  % VSS rail
  \draw[BookTeal, line width=0.7pt] (1.0,-2.2) -- (7.0,-2.2);
  \node[BookTeal, font=\small\bfseries] at (4.0,-2.5) {VSS};
  % Left inverter: P1 (top) and N1 (bottom)
  \node[hw state, minimum width=1.1cm, minimum height=0.65cm] (P1) at (2.5,1.8) {P1 (PMOS)};
  \node[hw state, minimum width=1.1cm, minimum height=0.65cm] (N1) at (2.5,-0.8) {N1 (NMOS)};
  % Right inverter: P2 (top) and N2 (bottom)
  \node[hw state, minimum width=1.1cm, minimum height=0.65cm] (P2) at (5.5,1.8) {P2 (PMOS)};
  \node[hw state, minimum width=1.1cm, minimum height=0.65cm] (N2) at (5.5,-0.8) {N2 (NMOS)};
  % Q node (left inverter output)
  \draw[BookAccent, line width=0.6pt] (P1.south) -- (N1.north);
  \node[BookAccent, font=\small\bfseries] at (1.8,0.5) {Q};
  % Qbar node (right inverter output)
  \draw[BookAccent, line width=0.6pt] (P2.south) -- (N2.north);
  \node[BookAccent, font=\small\bfseries] at (6.2,0.5) {$\overline{Q}$};
  % Cross-coupling: Q -> P2/N2 gate
  \draw[BookRed, line width=0.55pt] (2.5,0.5) -- (2.5,0.5) -| (4.6,1.8);
  \draw[BookRed, line width=0.55pt] (2.5,0.5) -| (4.6,-0.8);
  % Cross-coupling: Qbar -> P1/N1 gate
  \draw[BookRed, line width=0.55pt] (5.5,0.5) -| (3.4,1.8);
  \draw[BookRed, line width=0.55pt] (5.5,0.5) -| (3.4,-0.8);
  % P1 to VDD
  \draw[BookTeal, line width=0.5pt] (P1.north) -- ++(0,0.9);
  % P2 to VDD
  \draw[BookTeal, line width=0.5pt] (P2.north) -- ++(0,0.9);
  % N1 to VSS
  \draw[BookTeal, line width=0.5pt] (N1.south) -- ++(0,-0.9);
  % N2 to VSS
  \draw[BookTeal, line width=0.5pt] (N2.south) -- ++(0,-0.9);
  % Access transistors M5 and M6
  \node[hw compute, minimum width=0.9cm, minimum height=0.6cm] (M5) at (0.0,0.5) {M5};
  \node[hw compute, minimum width=0.9cm, minimum height=0.6cm] (M6) at (8.0,0.5) {M6};
  % M5 connections: BL -- M5 -- Q
  \draw[BookAccent, line width=0.6pt] (M5.east) -- (2.5,0.5);
  \draw[hw bus] (M5.west) -- ++(-1.2,0) node[left, hw label, font=\small\bfseries] {BL};
  % M6 connections: Qbar -- M6 -- BLbar
  \draw[BookAccent, line width=0.6pt] (M6.west) -- (5.5,0.5);
  \draw[hw bus] (M6.east) -- ++(1.2,0) node[right, hw label, font=\small\bfseries] {$\overline{BL}$};
  % Word line WL：只接 M5/M6 栅极，从顶部通道走线，不穿 P1/P2 框
  \draw[BookAmber, line width=0.8pt] (M5.north) -- (0.0,3.85);
  \draw[BookAmber, line width=0.8pt] (M6.north) -- (8.0,3.85);
  \draw[BookAmber, line width=0.5pt, dashed] (0.0,3.85) -- (8.0,3.85);
  \node[BookAmber, hw label, font=\small\bfseries] at (4.0,4.15) {WL};
  % Labels
  \node[hw label, text width=10cm] at (4.0,-3.4) {6T SRAM 单元完整逻辑图。红色线为交叉耦合反馈路径（Q 驱动 P2/N2 栅极，$\overline{Q}$ 驱动 P1/N1 栅极）。WL 同时控制 M5 和 M6 的栅极。BL/$\overline{BL}$ 是差分位线对。};
\end{pdfdiagram}
\diagramnote{6T SRAM 单元。P1/N1 和 P2/N2 交叉耦合形成双稳态锁存；M5/M6 是 NMOS 访问管，栅极共接字线 WL。WL=0 时 M5/M6 截止，锁存器自保持；WL=1 时导通，BL/$\overline{BL}$ 可读出或写入 Q/$\overline{Q}$。}

\hwkey{设计尺寸比}是 SRAM 单元设计的核心参数：

\begin{longtable}{L{0.28\linewidth}L{0.22\linewidth}L{0.40\linewidth}}
\toprule
参数 & 定义 & 典型值（28nm） \\
\midrule
Cell Ratio (CR) & $W_{N1}/W_{M5}$（下拉管/访问管） & 1.2--1.5 \\\tablerowrule
Pull-up Ratio (PR) & $W_{M5}/W_{P1}$（访问管/上拉管） & 1.0--1.5 \\\tablerowrule
驱动管宽长比 & $W/L$ of N1 & 最小尺寸（如 28nm: $W$=40nm） \\\tablerowrule
访问管宽长比 & $W/L$ of M5 & 略小于 N1 \\
\bottomrule
\end{longtable}

两个比值都以访问管 M5 为参照，名字也由此而来。\term{Cell Ratio}{CR}（单元比）衡量下拉管 N1 相对 M5 的强度，决定读操作时 Q=0 的节点能否顶住位线注入的电流而不被抬高；\term{Pull-up Ratio}{PR}（上拉比）衡量 M5 相对上拉管 P1 的强度，决定写操作时 M5 能否从 P1 手中把 Q 拉低。一句话：CR 管读、PR 管写——1.3 和 1.4 两节的定量分析将反复用到这两个比值，之后各节提到 CR/PR 时不再重复定义。

%======================================================================
\topic{保持操作：反馈环稳定性与漏电路径}
%======================================================================

WL=0 时，M5/M6 截止，6T 单元退化为一个\hwkey{交叉耦合锁存器}。假设 Q=1（VDD）、$\overline{Q}$=0（VSS）：

\begin{itemize}
  \item $\overline{Q}$=0 加在 P1/N1 的栅极上：P1 导通（PMOS 栅极低电平导通）、N1 截止，于是 Q 被 P1 拉在 VDD——与假设一致；
  \item Q=1 加在 P2/N2 的栅极上：N2 导通（NMOS 栅极高电平导通）、P2 截止，于是 $\overline{Q}$ 被 N2 拉在 VSS——同样与假设一致。
\end{itemize}

两条推理互为因果、彼此自洽，说明 Q=1/$\overline{Q}$=0 确实是一个稳定状态（由对称性，Q=0/$\overline{Q}$=1 是另一个）。这是一个\hwkey{正反馈环}：任何使 Q 略微下降的扰动，都会使 N2 略微关断、P2 略微导通，把 $\overline{Q}$ 拉高，进而使 N1 略微导通、P1 略微关断，把 Q 进一步拉低——直到翻转。但只要扰动幅度小于\hwkey{静态噪声裕量}（static noise margin, SNM，1.5 节给出定量定义与测量方法），反馈会把工作点推回稳态。

\hwkey{漏电路径}（保持状态下的非理想因素）：

\begin{longtable}{L{0.26\linewidth}L{0.24\linewidth}L{0.40\linewidth}}
\toprule
漏电来源 & 典型量级（28nm, 0.9V） & 影响 \\
\midrule
亚阈值漏电（N1/N2 截止时） & 1--10 pA/管 & 存储节点缓慢漂移 \\\tablerowrule
栅极隧穿漏电 & 0.1--1 pA/管 & 先进节点（$\leq$7nm）显著 \\\tablerowrule
M5/M6 截止漏电 & 1--5 pA/管 & 通过 BL 路径泄漏 \\\tablerowrule
PN 结反向漏电 & 0.5--2 pA/结 & 源/漏扩散区到衬底 \\
\bottomrule
\end{longtable}

四条路径中占主导的是前两条。\term{亚阈值漏电}{subthreshold leakage}指栅压低于阈值后沟道电流并不归零、而是随 $V_{GS}$ 指数减小的残余电流，它是深亚微米工艺保持功耗的主要来源；\term{栅极隧穿漏电}{gate tunneling leakage}指载流子以量子隧穿方式直接穿过栅氧化层的电流，氧化层每薄一分，隧穿概率按指数上升，因此只在先进节点变得显著。

在 28nm 工艺、0.9V 电源下，一个 6T 单元的总保持漏电约 \hwevidence{5--50 pA}。对于 1 Mbit 阵列，总静态漏电约 5--50 $\mu$A。在 7nm 及以下，栅极隧穿漏电占比增大，单元漏电可达 \hwevidence{50--200 pA}。

保持状态下的\hwkey{最小稳定电压}（data retention voltage, DRV）指数据尚不丢失的最低电源电压。它的物理来源是环路增益：VDD 降低后晶体管进入亚阈值区，跨导随电流指数下降，当两个反相器的环路增益不足 1 时，正反馈无力再放大扰动、对抗漏电，双稳态退化为单稳态，数据丢失。典型 DRV 在 28nm 为 \hwevidence{0.4--0.6V}，在 7nm 为 \hwevidence{0.3--0.5V}。DRV 决定了休眠模式下 SRAM 能把电压降到多低而不丢数据，1.11 节的 retention mode 正是利用这一点省电。

%======================================================================
\topic{读操作：电荷分享与灵敏放大（定量）}
%======================================================================

本小节把一次读操作拆成四个阶段，回答两个问题：位线上那几十毫伏的差分信号从哪里来，以及它怎样被放大成干净的 0/1。读之前必须先\term{预充}{precharge}：用预充管把一对位线都充到 VDD，然后关断预充管，使位线处于“高电平但浮空”的待发状态——之后发生的一切电压变化都以这个起点为参照。读操作分四个阶段：

\begin{lstlisting}
// 注释（推进）：逐行追踪发起者、所有权和可见性；请求被接受不等于事务完成。
// 注释（边界）：以采样沿、状态位或 completion 为准，再允许消费者使用结果。
READ OPERATION (step by step):
  Phase 1 - Precharge:
    BL = VDD, BLbar = VDD    // 位线对预充到满高
    WL = 0                    // 字线关闭，单元隔离
  Phase 2 - Wordline assert:
    WL = VDD                  // M5, M6 导通
    // 若 Q=0: BL 通过 M5 -> N1 向 VSS 放电
    // 若 Q=1: BL 保持 VDD（P1 上拉 vs M5 无放电路径）
  Phase 3 - Charge sharing / differential develop:
    // BL 下降 delta_V，BLbar 保持（或微升）
    // 差分电压 V(BL) - V(BLbar) 逐渐增大
  Phase 4 - Sense amplifier latch:
    // 交叉耦合感放检测差分，正反馈放大到全摆幅
    // 输出 DOUT = 0 (if Q=0) or 1 (if Q=1)
\end{lstlisting}

\hwkey{定量分析}——电荷分享模型：

设存储节点 Q=0（即 Q 端电压 $\approx$ 0V），BL 预充到 VDD=0.9V。WL 打开后，BL 通过 M5 和 N1 向地放电。但放电不是到 0V——因为 N1 的驱动能力有限，且 BL 电容远大于单元内部电容。实际放电速度由\hwkey{放电通路}决定：读 Q=0 侧时 P1 栅接 $\overline{Q}$=VDD 处于截止，BL 的放电电流只能经 M5、N1 串联流向地，放电极限由 N1 的下拉能力决定。

简化电荷分享模型（读 Q=0 时 P1 截止，暂不考虑其影响）：

$$\Delta V_{BL} = \frac{C_{cell}}{C_{cell} + C_{BL}} \times V_{DD}$$

其中 $C_{cell}$ 是存储节点对地电容（含 M5 漏极、N1 漏极、P1 漏极的扩散电容和连线电容），$C_{BL}$ 是位线总电容。

\begin{longtable}{L{0.30\linewidth}L{0.22\linewidth}L{0.38\linewidth}}
\toprule
参数 & 典型值（28nm） & 来源 \\
\midrule
$C_{cell}$（存储节点电容） & 1--3 fF & 扩散结电容 + 连线 \\\tablerowrule
$C_{BL}$（位线电容） & 200--400 fF & 128--256 行单元的 M5 漏极 + 金属线 \\\tablerowrule
$C_{cell}/C_{BL}$ & 1:100 -- 1:200 & 比值极小 \\\tablerowrule
$\Delta V_{BL}$（单次电荷分享） & 5--15 mV & 理论初始值 \\\tablerowrule
实际读放电（含电流竞争） & 50--100 mV & 感放触发前累积放电 \\
\bottomrule
\end{longtable}

表中最后两行描述的是同一条 BL 电压曲线的两段，不要误读为两个独立过程。\term{电荷分享}{charge sharing}只是 WL 刚打开瞬间的电压阶跃：预充到 VDD 的 BL 与接近 0V 的存储节点突然连通，两边电荷按电容比例重新分配，在皮秒量级内形成 5--15 mV 的初始落差。阶跃之后放电并未停止——N1 持续导通，以近似恒定的电流把 BL 上的电荷泄向地，BL 电压随时间近似线性下降，速率约为 $I_{N1}/C_{BL}$，按 28nm 典型值估算在每纳秒几十到一百多 mV 的量级。表中“实际读放电 50--100 mV”是感放触发时刻（WL 打开后约 0.3--0.5 ns）采样到的累积差分，而不是另一次电荷分享。设计中调节感放触发时刻，就是在“多等一会儿让差分长大”与“尽早读出求快”之间取折中。

为什么不让 BL 自行放电到全摆幅再判断？因为几百 fF 的 BL 电容靠单元里最小尺寸的晶体管放电需要数纳秒，速度无法接受。解决办法是用\term{灵敏放大器}{sense amplifier, SA}：一个由时钟使能的交叉耦合锁存器（核心结构与 6T 单元相同），使能瞬间把毫伏级差分接入正反馈环，在亚纳秒内放大到全摆幅。这样“等电容放完电”就变成了“等电压差长大到可判别”，这是 SRAM 与 DRAM 共用的加速思路。感放并非理想器件：两侧晶体管的阈值失配等效为一个输入\term{失调电压}{offset voltage}，典型小于 20 mV，差分信号必须明显大于失调电压，读出才可靠。

\hwevidence{关键数据}：在 28nm 工艺、VDD=0.9V、256 行阵列中，$C_{BL} \approx$ 300 fF，读放电幅度约 \hwevidence{60--80 mV}，感放触发延迟约 \hwevidence{0.3--0.5 ns}，总读访问时间（WL 有效到 DOUT 稳定）约 \hwevidence{0.8--1.5 ns}。

\begin{pdfdiagram}[x=1cm,y=1cm]
  % Time axis
  \draw[BookRule, line width=0.5pt, ->] (0,0) -- (11,0) node[right, hw label] {时间 (ns)};
  % WL waveform
  \draw[BookAmber, line width=0.7pt] (0.5,3.8) -- (1.5,3.8) -- (1.5,4.4) -- (9.5,4.4) -- (9.5,3.8) -- (10.5,3.8);
  \node[hw label, anchor=east] at (0.4,4.1) {WL};
  % BL waveform (discharging)
  \draw[BookTeal, line width=0.8pt] (0.5,2.8) -- (1.5,2.8) -- (2.0,2.75) -- (3.0,2.6) -- (4.0,2.5) -- (5.0,2.45) -- (5.5,2.4) -- (6.0,1.2) -- (10.5,1.2);
  \node[hw label, anchor=east] at (0.4,2.8) {BL (Q=0)};
  % BLbar waveform (stays high)
  \draw[BookAccent, line width=0.8pt] (0.5,2.2) -- (1.5,2.2) -- (5.5,2.2) -- (6.0,3.2) -- (10.5,3.2);
  \node[hw label, anchor=east] at (0.4,2.2) {$\overline{BL}$ (Q=1)};
  % Annotations
  \draw[BookRule, line width=0.3pt, dashed] (1.5,4.6) -- (1.5,-0.3);
  \draw[BookRule, line width=0.3pt, dashed] (5.5,4.6) -- (5.5,-0.3);
  \node[hw label] at (1.0,4.8) {预充};
  \node[hw label] at (3.5,4.8) {差分发展};
  \node[hw label] at (8.0,4.8) {感放锁存};
  % Delta V annotation
  \draw[BookRed, line width=0.5pt, <->] (5.7,2.45) -- (5.7,2.2);
  \node[BookRed, font=\tiny\bfseries] at (6.3,2.32) {$\Delta V$=60mV};
  % Sense amp trigger
  \node[hw label, font=\tiny] at (5.5,-0.5) {感放触发点};
  \node[hw label, text width=9cm] at (5.5,-1.2) {读操作时序：WL 打开后 BL 缓慢放电（Q=0 侧），$\overline{BL}$ 保持高（Q=1 侧）。差分电压超过感放的最小可辨阈值（灵敏度下限，约 50--80mV，远大于失调电压）后，正反馈在 0.3ns 内放大到全摆幅。};
\end{pdfdiagram}
\diagramnote{SRAM 读操作波形。注意读操作是\hwkey{非破坏性}的——读 Q=0 侧时 P1 栅接 $\overline{Q}$=VDD 处于截止，只要 N1 相对 M5 足够强（CR $>$ 1），Q 就不会被抬过翻转点，读完后单元数据不变。}

\hwkey{读操作的关键约束}：读过程中，预充到 VDD 的 BL 经 M5 向 Q=0 的节点注入电流，Q 被向上抬；同时 N1 把 Q 拉向地——Q 的稳态电压由 M5 与 N1 的分压决定。如果 M5 太强（相对于 N1），Q 可能被抬高到 N2 的阈值电压以上，导致单元意外翻转——这就是\hwkey{读扰动（read disturb）}。因此 N1 必须比 M5 强（即 CR = $W_{N1}/W_{M5}$ $>$ 1），保证读稳定性。

%======================================================================
\topic{写操作：驱动能力竞争与 CR/PR}
%======================================================================

“从 bit 到字到阵列到系统：结构递进”一节的读操作刻意不去扰动单元，本小节的写操作则相反：目标是\hwkey{强行翻转}锁存器状态。假设当前 Q=1、$\overline{Q}$=0，要写入 0：

\begin{lstlisting}
// 注释（推进）：逐行追踪发起者、所有权和可见性；请求被接受不等于事务完成。
// 注释（边界）：以采样沿、状态位或 completion 为准，再允许消费者使用结果。
WRITE OPERATION (writing 0 to cell with Q=1, Qbar=0):
  Phase 1 - Drive bitlines:
    BL = 0 (VSS)             // 写驱动器拉低 BL
    BLbar = VDD              // 写驱动器拉高 BLbar
    WL = 0                    // 字线尚未打开
  Phase 2 - Assert wordline:
    WL = VDD                  // M5, M6 导通
    // BL=0 通过 M5 拉低 Q（对抗 P1 上拉）
    // BLbar=VDD 通过 M6 拉高 Qbar（对抗 N2 下拉）
  Phase 3 - Cell flip:
    // Q 被拉低到 N2 阈值以下 -> N2 截止
    // Qbar 被拉高到 P1 阈值以上 -> P1 截止
    // 正反馈接管：N1 导通拉 Q->0, P2 截止让 Qbar->VDD
    // 锁存器翻转到新稳态：Q=0, Qbar=1
  Phase 4 - Deassert wordline:
    WL = 0                    // 新数据锁存
    BL, BLbar 恢复预充
\end{lstlisting}

\hwkey{写操作的核心矛盾}：写驱动器通过 M5 拉低 Q，但 P1（PMOS 上拉管）同时把 Q 拉向 VDD。写入成功的条件是：

$$I_{M5}(\text{pull-down}) > I_{P1}(\text{pull-up}) \quad \text{at } V_Q = V_{th,N2}$$

即 M5 在 Q 被拉到 N2 阈值电压（约 0.3--0.4V）时提供的下拉电流，必须大于 P1 在此电压下的上拉电流。

注意写入的主攻方向是\hwkey{拉低} Q，而不是抬高 $\overline{Q}$：NMOS 访问管传 0 没有电平损失，传 1 却要损失约一个阈值电压（栅压 VDD 减去 $V_{th}$ 后管子已近截止），靠 M6 把 $\overline{Q}$ 推向 VDD 既慢又不可靠。因此写电路总是由写驱动器把对应一侧的位线拉到 0，让另一侧靠交叉耦合反馈自己跟上。

1.1 节定义的两个尺寸比在这里各就其位。写能力由 PR = $W_{M5}/W_{P1}$ 决定：PR 大意味着 M5 相对 P1 强，写 0 时容易把 Q 拉过翻转点。读稳定性则由 CR = $W_{N1}/W_{M5}$ 决定：读 Q=0 侧时 P1 的栅极接在 $\overline{Q}$=VDD 上处于截止，与 PR 无关，Q 能否稳住只看 N1 是否比 M5 强。一个比值管写、一个管读，这正是读写裕量可以在一定程度上独立调节的结构基础。

典型设计值：CR = 1.2--1.5，PR = 1.0--1.5。在先进节点（7nm/5nm），由于 FinFET 的量化宽度（fin 数为整数），CR 和 PR 的调节粒度变粗，设计空间收窄。

\begin{pdfdiagram}[x=1cm,y=1cm]
  % Write operation: current competition
  \node[hw state, minimum width=1.0cm, minimum height=0.6cm] (P1w) at (2.0,2.0) {P1};
  \node[hw compute, minimum width=1.0cm, minimum height=0.6cm] (M5w) at (2.0,0.0) {M5};
  \node[hw state, minimum width=1.0cm, minimum height=0.6cm] (N1w) at (2.0,-1.5) {N1};
  % Connections
  \draw[BookTeal, line width=0.6pt] (P1w.north) -- ++(0,0.6) node[above, hw label] {VDD};
  \draw[BookAccent, line width=0.6pt] (P1w.south) -- (M5w.north);
  \node[BookAccent, font=\small\bfseries] at (2.8,1.0) {Q};
  \draw[BookAccent, line width=0.6pt] (M5w.south) -- (0.4,-0.3) node[left, hw label, font=\small\bfseries] {BL=0};
  \draw[BookTeal, line width=0.5pt] (N1w.south) -- ++(0,-0.5) node[below, hw label] {VSS};
  % Current arrows
  \draw[BookRed, line width=0.7pt, ->] (1.2,2.0) -- (1.2,1.0) node[midway, left, font=\tiny, BookRed] {$I_{P1}$};
  \draw[BookTeal, line width=0.7pt, ->] (1.2,0.0) -- (1.2,-0.8) node[midway, left, font=\tiny, BookTeal] {$I_{M5}$};
  % Condition
  \node[hw label, text width=7cm] at (5.5,0.5) {写入条件：$I_{M5} > I_{P1}$ 在 $V_Q = V_{th}$ 时。\\写驱动器（大尺寸 NMOS）通过 BL=0 和 M5 拉低 Q，\\必须克服 P1 的上拉。};
  % Write driver：正交走线，方向 Write Driver -> BL -> M5
  \node[hw compute, minimum width=1.4cm, minimum height=0.7cm] (wd) at (0.4,-1.6) {Write Driver};
  \draw[hw bus] (wd.north) -- (0.4,-0.3);
  \node[hw label, font=\tiny] at (0.4,-2.35) {大尺寸 NMOS 下拉};
\end{pdfdiagram}
\diagramnote{写操作中的电流竞争。P1 试图把 Q 拉回 VDD，M5+写驱动器试图把 Q 拉到 0。写入成功要求 M5 路径的驱动能力超过 P1。}

%======================================================================
\topic{静态噪声裕量（SNM）与蝶形图}
%======================================================================

1.2 节定性地说过“扰动不超过静态噪声裕量就翻不了”，本小节给出这个裕量的严格定义与工程上通用的测量方法。\hwkey{静态噪声裕量（Static Noise Margin, SNM）}是衡量 SRAM 单元在保持状态下抗干扰能力的标准度量。测量方法：

\begin{enumerate}
  \item 画出左反相器（P1/N1）的电压传输特性（VTC）：$V_{out,Q} = f(V_{in,\overline{Q}})$
  \item 画出右反相器（P2/N2）的 VTC：$V_{out,\overline{Q}} = g(V_{in,Q})$
  \item 将第二条 VTC 以 $V_{out}$ 轴为对称轴镜像（即交换横纵坐标）
  \item 两条曲线叠加形成\hwkey{蝶形图（butterfly curve）}，围成两个“眼”
  \item 每个眼中能放入的\hwkey{最大内切正方形}的边长即为 SNM（取两个眼中较小者）
\end{enumerate}

\begin{classicnote}
为什么取最大内切正方形？把两个反相器首尾相接成环路后，外部干扰可以等效为同时串联在两个输入端的等幅噪声源；两条 VTC 曲线之间能容纳的最大正方形边长，正是环路在保持双稳态的前提下、两个节点可同时承受的最大串联噪声幅度。这个“最大正方形”判据由 Seevinck 等人于 1987 年提出，至今仍是 SRAM 文献的标准度量。
\end{classicnote}

\begin{pdfdiagram}[x=1cm,y=1cm]
  % Axes
  \draw[BookRule, line width=0.5pt, ->] (0,0) -- (6.0,0) node[right, hw label] {$V_Q$ (V)};
  \draw[BookRule, line width=0.5pt, ->] (0,0) -- (0,6.0) node[above, hw label] {$V_{\overline{Q}}$ (V)};
  % VDD labels
  \node[hw label] at (5.5,-0.3) {0.9};
  \node[hw label] at (-0.3,5.5) {0.9};
  % VTC curve 1 (inverter 1: input Qbar, output Q) - normal S shape（红）
  \draw[BookRed, line width=0.9pt] (0.3,5.2) .. controls (1.2,5.2) and (1.8,4.8) .. (2.5,2.8) .. controls (3.2,0.8) and (3.8,0.4) .. (5.2,0.4);
  % VTC curve 2 (inverter 2 mirrored) - inverted S（蓝）
  \draw[BookAccent, line width=0.9pt] (5.2,5.2) .. controls (3.8,5.2) and (3.2,4.8) .. (2.5,2.8) .. controls (1.8,0.8) and (1.2,0.4) .. (0.3,0.4);
  % SNM squares (two eyes)
  \draw[BookAmber, line width=0.8pt] (0.8,3.2) rectangle (2.2,4.6);
  \draw[BookAmber, line width=0.8pt] (3.3,1.0) rectangle (4.7,2.4);
  % SNM label
  \node[BookAmber, font=\small\bfseries] at (1.5,3.9) {SNM$_1$};
  \node[BookAmber, font=\small\bfseries] at (4.0,1.7) {SNM$_2$};
  % Stable points
  \fill[BookRed] (0.5,5.0) circle (0.08);
  \fill[BookRed] (5.0,0.5) circle (0.08);
  \node[BookRed, font=\tiny\bfseries] at (0.9,5.3) {稳态 1};
  \node[BookRed, font=\tiny\bfseries] at (5.3,0.9) {稳态 2};
  % Meta-stable point
  \fill[BookMuted] (2.5,2.8) circle (0.06);
  \node[hw label, font=\tiny] at (3.0,3.1) {亚稳态};
  % Caption
  \node[hw label, text width=9cm] at (3.0,-1.0) {蝶形图：红色曲线为反相器 1 的 VTC，蓝色曲线为反相器 2 的镜像 VTC。两个“眼”中的最大内切正方形边长 = SNM。SNM = min(SNM$_1$, SNM$_2$)。工艺波动使两个眼不对称，较小者决定单元裕量。};
\end{pdfdiagram}
\diagramnote{SRAM 蝶形图与 SNM 测量。两个稳态点（红点）位于对角；亚稳态点（灰点）在中心。SNM 越大，单元越不容易被噪声翻转。}

\hwkey{不同工艺节点的 SNM 典型值}：

\begin{longtable}{L{0.12\linewidth}L{0.12\linewidth}L{0.18\linewidth}L{0.18\linewidth}L{0.30\linewidth}}
\toprule
工艺节点 & VDD (V) & 标称 SNM (mV) & 3$\sigma$ 最差 SNM & 备注 \\
\midrule
65 nm & 1.0 & 250--350 & 180--250 & 平面 MOSFET，裕量充足 \\\tablerowrule
28 nm & 0.9 & 150--250 & 100--160 & 平面 MOSFET 最后节点 \\\tablerowrule
14 nm & 0.8 & 120--180 & 70--120 & FinFET 引入，波动减小 \\\tablerowrule
7 nm & 0.7 & 80--150 & 40--90 & EUV 光刻，fin 数量少 \\\tablerowrule
5 nm & 0.65 & 60--120 & 30--70 & 需要 assist 电路 \\\tablerowrule
3 nm & 0.6 & 50--100 & 20--50 & GAA 纳米片，设计极困难 \\
\bottomrule
\end{longtable}

表中“3$\sigma$ 最差 SNM”指计入工艺随机波动后、比标称值差 3 个标准差的单元的 SNM。对百万级单元的阵列而言，起决定作用的是分布尾部的最差单元而不是标称值，因此设计必须按 3$\sigma$（高良率要求时甚至 6$\sigma$）列来校核。

\hwkey{SNM 缩小的原因}：
\begin{enumerate}
  \item \hwkey{VDD 降低}：SNM $\propto$ VDD（近似线性），VDD 从 1.0V 降到 0.65V，SNM 直接缩小 35\%
  \item \hwkey{阈值电压波动}（$V_{th}$ mismatch）：$\sigma_{V_{th}} \propto 1/\sqrt{W \cdot L}$，面积缩小后波动增大
  \item \hwkey{短沟道效应}：DIBL（漏致势垒降低）使亚阈值斜率退化，反相器增益下降
  \item \hwkey{FinFET 量化}：fin 数只能取整数（1, 2, 3...），无法精细调节驱动比
\end{enumerate}

%======================================================================
\topic{读稳定性 vs 写能力的权衡}
%======================================================================

前面三节分别考察了保持、读、写，本小节把它们放在一起：加大同一个晶体管，往往一边受益、一边受损，这是 SRAM 设计中\hwkey{最根本的矛盾}。注意本节的噪声裕量与 1.5 节测的不是同一个：1.5 节的蝶形图对应 WL=0 的\emph{保持}状态，而读操作时 WL=1、位线预充在 VDD，M5 持续向 Q=0 节点注入电流，等效于在随机噪声之上再叠加一个固定方向的扰动，因此读状态的噪声裕量总是显著小于保持状态——读才是单元最容易翻转的时刻。

\begin{longtable}{L{0.18\linewidth}L{0.36\linewidth}L{0.36\linewidth}}
\toprule
设计选择 & 对读稳定性的影响 & 对写能力的影响 \\
\midrule
M5 更强（$W_{M5}$ 增大） & 读时 M5 把 BL 电压传到 Q，更容易把 Q 拉高 → \hwrisk{读稳定性下降} & 写时 M5 提供更多下拉电流 → \hwkey{写能力增强} \\\tablerowrule
P1 更强（$W_{P1}$ 增大） & 读时 P1 把 Q 牢牢拉在 VDD → \hwkey{读稳定性增强} & 写时 P1 对抗 M5 更强 → \hwrisk{写能力下降} \\\tablerowrule
N1 更强（$W_{N1}$ 增大） & 读时 N1 提供强放电路径 → 读信号大 & 写时 N1 帮助拉低 Q → 写容易 \\
\bottomrule
\end{longtable}

\hwkey{量化权衡}：定义\hwkey{读 SNM}（read SNM, RSNM）为 WL=1、位线预充到 VDD 时测得的 SNM，即读状态下的蝶形图内切正方形边长，它是所有工作状态中最小的噪声裕量。\hwkey{写裕量}（write margin, WM）沿用文献中的常用定义：写 0 时把 BL 从 VDD 开始逐渐下拉，单元恰好翻转那一瞬的 $V_{DD}-V_{BL}$ 即为 WM。WM 越大，说明 BL 不必拉多低单元就已翻转，写入越容易、裕量越大——与 RSNM 一样，数值大才是好事。

\begin{longtable}{L{0.22\linewidth}L{0.26\linewidth}L{0.26\linewidth}L{0.16\linewidth}}
\toprule
工艺 & RSNM (mV) & WM (mV) & 设计难点 \\
\midrule
65 nm & 200--300 & 400--500 & 裕量充足 \\\tablerowrule
28 nm & 120--200 & 250--350 & 开始紧张 \\\tablerowrule
7 nm & 60--120 & 150--250 & 需要 assist \\\tablerowrule
5 nm & 40--90 & 100--180 & 强依赖 assist \\
\bottomrule
\end{longtable}

对照表中趋势：工艺每前进一代，RSNM 与 WM 同步收窄。在 7nm 及以下的最差\hwkey{工艺角}（process corner，即掺杂浓度、器件尺寸、温度、电压波动的最坏组合）下，两者都可能退化到接近于零——同一颗芯片上可能既有读时会翻转的单元，又有写不进去的单元。这就是\hwkey{写辅助电路}（write assist）和\hwkey{读辅助电路}（read assist）存在的根本原因，也是下一小节的主题。

%======================================================================
\topic{写辅助电路（Write/Read Assist）}
%======================================================================

当标称设计无法满足所有工艺角的读写裕量时，引入辅助电路在关键时刻临时改变偏置条件：

\begin{longtable}{L{0.18\linewidth}L{0.32\linewidth}L{0.20\linewidth}L{0.20\linewidth}}
\toprule
辅助技术 & 原理 & 裕量改善 & 代价 \\
\midrule
\hwkey{负位线}（Negative BL） & 写时把 BL 驱动到 $-$0.1 $\sim$ $-$0.2V（低于 VSS），增大 M5 的 $V_{GS}$ & 写裕量 +80--150 mV & 需要负电压发生器；增加面积 3--5\% \\\tablerowrule
\hwkey{字线下压}（WL underdrive） & 读时 WL 不拉到满 VDD，而是 0.7$\times$VDD，减弱 M5 导通 & 读 SNM +30--60 mV & 读速度变慢 +0.1--0.2 ns \\\tablerowrule
\hwkey{VDD 下压}（VDD collapse） & 写时临时把单元的 VDD 拉到 0.5$\times$VDD，削弱 P1 上拉 & 写裕量 +100--200 mV & 需要 per-bank VDD 开关；恢复时间 \\\tablerowrule
\hwkey{读 VDD 提升}（Read VDD boost） & 读时把单元 VDD 提升到 1.1$\times$VDD，增强 P1 上拉 & 读 SNM +40--80 mV & 需要升压电路；可靠性应力 \\\tablerowrule
\hwkey{体偏置}（Body bias） & 通过衬底偏置调节 $V_{th}$：写时正向体偏（降低 NMOS $V_{th}$） & 读/写均可调 & 需要 triple-well 隔离；面积 +10\% \\
\bottomrule
\end{longtable}

把这五种技术放在一起看，原理是统一的：读写困难的根源都是有效 CR/PR 不够大，辅助电路只是在需要的那几十皮秒里临时改变偏置条件——写时让 M5 显得更强或 P1 显得更弱（等效于增大 PR），读时让 M5 显得更弱或 P1 显得更强（等效于增大 CR）。之所以只在操作瞬间施加而不永久改变偏置，是因为这些偏置都偏离正常设计点：负压与过压会加剧器件老化，VDD 下压会侵蚀保持裕量，长期施加得不偿失。

\begin{pdfdiagram}[x=1cm,y=1cm]
  % Negative BL diagram
  \node[hw label, font=\small\bfseries] at (2.0,3.5) {负位线写辅助};
  \draw[BookRule, line width=0.5pt, ->] (0,0) -- (5,0) node[right, hw label] {时间};
  \draw[BookRule, line width=0.5pt, ->] (0,0) -- (0,3.2) node[above, hw label] {电压};
  % Normal BL=0
  \draw[BookTeal, line width=0.7pt] (0.5,0.5) -- (4.5,0.5);
  \node[hw label, anchor=east] at (0.4,0.5) {0V};
  % Negative BL
  \draw[BookRed, line width=0.7pt] (0.5,0.5) -- (1.5,0.5) -- (1.5,-0.3) -- (3.5,-0.3) -- (3.5,0.5) -- (4.5,0.5);
  \node[BookRed, font=\tiny\bfseries] at (2.5,-0.6) {BL = $-$0.15V};
  % VDD level
  \draw[BookAccent, line width=0.5pt, dashed] (0.5,2.8) -- (4.5,2.8);
  \node[hw label, anchor=east] at (0.4,2.8) {VDD};
  % VGS annotation
  \draw[BookAmber, line width=0.5pt, <->] (4.0,2.8) -- (4.0,-0.3);
  \node[BookAmber, font=\tiny\bfseries, rotate=90] at (4.4,1.2) {$V_{GS}$ 增大};
  % VDD collapse diagram
  \node[hw label, font=\small\bfseries] at (8.0,3.5) {VDD 下压写辅助};
  \draw[BookRule, line width=0.5pt, ->] (6,0) -- (11,0) node[right, hw label] {时间};
  \draw[BookRule, line width=0.5pt, ->] (6,0) -- (6,3.2) node[above, hw label] {电压};
  % Normal VDD
  \draw[BookTeal, line width=0.5pt, dashed] (6.5,2.8) -- (10.5,2.8);
  \node[hw label, anchor=east] at (6.4,2.8) {VDD};
  % Collapsed VDD
  \draw[BookRed, line width=0.7pt] (6.5,2.8) -- (7.5,2.8) -- (7.5,1.4) -- (9.5,1.4) -- (9.5,2.8) -- (10.5,2.8);
  \node[BookRed, font=\tiny\bfseries] at (8.5,1.1) {VDD$_{cell}$ = 0.5$\times$VDD};
  % P1 weakened
  \node[hw label, font=\tiny] at (8.5,0.5) {P1 上拉电流减半 → 写容易};
\end{pdfdiagram}
\diagramnote{两种写辅助技术。左：负位线增大 M5 的 $V_{GS}$（从 VDD 变为 VDD+0.15V），增强下拉。右：VDD 下压削弱 P1 上拉，降低写翻转阈值。}

\hwkey{工业实践}：\hwevidence{TSMC 7nm SRAM 宏}使用负位线（$-$100mV）+ WL 下压的组合；\hwevidence{Samsung 5nm} 使用 VDD collapse + 体偏置；\hwevidence{Intel 10nm（Intel 7）}使用读 VDD boost + 负位线。

%======================================================================
\topic{从单元到阵列：外围电路}
%======================================================================

一个实用的 SRAM 宏（macro，即可作为模块例化的成品存储块）由\hwkey{存储阵列}和\hwkey{外围电路}组成。以 28nm 工艺、512 Kbit（64 KiB）SRAM 宏为例，阵列组织为 256 行 $\times$ 2048 列（含纠错码位；纠错码 ECC 在 1.9 与 1.12 节介绍）。

\begin{pdfdiagram}[x=1cm,y=1cm]
  % Array block
  \draw[BookRule, line width=0.6pt, fill=BookCodeBg] (2.0,1.0) rectangle (8.0,5.0);
  \node[hw label, font=\small\bfseries] at (5.0,3.0) {存储阵列\\256 行 $\times$ 2048 列};
  % Row decoder (left)
  \node[hw control, minimum width=1.2cm, minimum height=2.5cm] (rdec) at (0.5,3.0) {行译码器\\(Row Dec)};
  \draw[hw bus] (rdec.east) -- (2.0,3.0);
  % Wordlines
  \foreach \y in {1.5, 2.2, 2.9, 3.6, 4.3} {
    \draw[BookAmber, line width=0.4pt] (2.0,\y) -- (8.0,\y);
  }
  \node[hw label, font=\tiny, align=center] at (1.55,3.38) {WL\\{[255:0]}};
  % Column mux (bottom)
  \node[hw compute, minimum width=3.0cm, minimum height=0.7cm] (cmux) at (5.0,0.3) {列多路选择器 (8:1 Mux)};
  \draw[hw bus] (5.0,1.0) -- (5.0,0.65);
  % Sense amplifiers
  \node[hw state, minimum width=2.5cm, minimum height=0.6cm] (sa) at (5.0,-0.6) {灵敏放大器组 (SA)};
  \draw[hw bus] (cmux.south) -- (sa.north);
  % Write drivers
  \node[hw compute, minimum width=2.0cm, minimum height=0.6cm] (wd) at (9.5,0.3) {写驱动器};
  \draw[hw bus] (wd.west) -- (8.0,0.3) -- (8.0,1.0);
  % Column decoder
  \node[hw control, minimum width=1.2cm, minimum height=0.8cm] (cdec) at (9.5,-0.6) {列译码器};
  % 列译码输出驱动列 mux 的选择端（正交折线从 mux 东侧进入，不进 SA）
  \draw[hw bus] (cdec.west) -- (7.2,-0.6) -- (7.2,0.3) -- (cmux.east);
  % Address input
  \node[hw label, font=\small\bfseries] at (0.5,5.5) {ADDR[15:0]};
  \draw[hw bus] (0.5,5.3) -- (rdec.north);
  % 列地址从顶部通道绕到列译码器（不进阵列框）
  \draw[hw bus] (0.5,5.2) -- (10.9,5.2) -- (10.9,-0.6) -- (cdec.east);
  % Data I/O
  \node[hw label, font=\small\bfseries] at (5.0,-1.5) {DIN / DOUT};
  \draw[hw bus] (sa.south) -- (5.0,-1.3);
  % Bitlines annotation
  \foreach \x in {2.5, 3.5, 4.5, 5.5, 6.5, 7.5} {
    \draw[BookTeal, line width=0.3pt] (\x,1.0) -- (\x,5.0);
  }
  \node[hw label, font=\tiny, rotate=90] at (8.3,3.0) {BL pairs};
  % Caption
  \node[hw label, text width=9cm] at (5.0,-2.3) {SRAM 宏结构：行译码器选中一条 WL；列译码器通过列 mux 选中一组 BL 对送到灵敏放大器；写驱动器在写操作时驱动 BL 对。};
\end{pdfdiagram}
\diagramnote{SRAM 宏内部结构。行地址选择 WL，列地址通过列 mux 选择 BL 组。灵敏放大器位于列 mux 和数据总线之间。}

\hwkey{各外围模块详解}：

\begin{longtable}{L{0.16\linewidth}L{0.38\linewidth}L{0.36\linewidth}}
\toprule
模块 & 功能 & 设计要点 \\
\midrule
\hwkey{行译码器} & 将行地址（如 8 bit → 256 条 WL）译码为一条 WL 有效 & 驱动长 WL（横跨整行）；WL 电容 $\propto$ 列数；需要大尺寸驱动管 \\\tablerowrule
\hwkey{字线驱动器} & 提供 WL 的充放电电流 & WL 延迟 = $R_{driver} \times C_{WL}$；横跨 2048 列的 WL 电容约 0.5--1 pF \\\tablerowrule
\hwkey{位线组织} & 每列一对 BL/$\overline{BL}$，贯穿所有行 & BL 电容 = 每行 M5 漏极电容 $\times$ 行数 + 金属线电容；256 行约 200--400 fF \\\tablerowrule
\hwkey{列多路选择器} & 从 N 条 BL 对中选 1 条送到 SA & 4:1 / 8:1 / 16:1；减少 SA 数量（面积）但增加 BL 负载 \\\tablerowrule
\hwkey{灵敏放大器} & 检测 BL 差分电压（50--100 mV）并放大到全摆幅 & 交叉耦合锁存型或电流镜型；offset $<$ 20 mV \\\tablerowrule
\hwkey{写驱动器} & 写时驱动 BL/$\overline{BL}$ 为互补电平 & 驱动能力 $>$ 单元 P1 上拉；通常是大尺寸反相器对 \\\tablerowrule
\hwkey{列译码器} & 将列地址译码为列 mux 的选择信号 & 与行译码器配合确定唯一被访问的单元 \\
\bottomrule
\end{longtable}

\hwkey{列 mux 比例的设计权衡}：mux 比例回答的问题是“每个灵敏放大器服务多少对位线”。比例越高，SA 数量越少、面积越省；但 mux 的传输管串在 BL 与 SA 之间，既让每条 BL 多挂一组开关的寄生电容，又给读出路径插入额外电阻，速度随之下降。注意它不改变位线的物理长度——BL 长度由行数决定，mux 比例决定的只是每个 SA 挂接的 BL 对数及其扇出负载。

\begin{longtable}{L{0.12\linewidth}L{0.16\linewidth}L{0.32\linewidth}L{0.30\linewidth}}
\toprule
Mux 比例 & SA 数量 & 每个 SA 挂接的 BL 对数与负载 & 适用场景 \\
\midrule
1:1 & 最多 & 1 对，直接相连，负载最小、速度最快 & L1 cache（速度优先） \\\tablerowrule
4:1 & 1/4 & 4 对，经一级 mux 选择 & L2 cache \\\tablerowrule
8:1 & 1/8 & 8 对，mux 树加深、插入延迟增大 & L2/L3（面积优先） \\\tablerowrule
16:1 & 1/16 & 16 对，mux 级数与扇出负载最大 & L3/大容量 SRAM \\
\bottomrule
\end{longtable}

%======================================================================
\topic{阵列组织：Cache 中的 SRAM}
%======================================================================

CPU 中的 SRAM 阵列根据用途有不同的组织方式：

\begin{longtable}{L{0.14\linewidth}L{0.14\linewidth}L{0.14\linewidth}L{0.14\linewidth}L{0.34\linewidth}}
\toprule
参数 & L1D & L1I & L2 & L3 (LLC) \\
\midrule
容量 & 32--64 KiB & 32--64 KiB & 256 KiB--2 MiB & 8--128 MiB（共享） \\\tablerowrule
端口数 & 2（读+写） & 1（读为主） & 1 或 2 & 1 \\\tablerowrule
访问延迟 & 3--5 周期 & 3--5 周期 & 8--14 周期 & 25--50 周期 \\\tablerowrule
Bank 数 & 2--4 & 1--2 & 4--8 & 8--32 \\\tablerowrule
行宽 & 64B + ECC & 64B + ECC & 64B + ECC & 64B + ECC \\
\bottomrule
\end{longtable}

\hwkey{单端口 vs 双端口}：

\begin{longtable}{L{0.18\linewidth}L{0.36\linewidth}L{0.36\linewidth}}
\toprule
类型 & 结构 & 应用 \\
\midrule
\hwkey{单端口}（1-port） & 1 对 BL/$\overline{BL}$，1 个 WL；同一周期只能读或写 & L1I、L2、L3（读写分时） \\\tablerowrule
\hwkey{双端口}（2-port） & 2 对 BL/$\overline{BL}$，2 个 WL；可同时一读一写 & 寄存器堆（RF）：读端口 $\times$2 + 写端口 $\times$1--2 \\\tablerowrule
\hwkey{多端口}（multi-port） & N 对 BL；面积 $\propto$ 端口数 & 极少用（面积太大）；寄存器堆通常用 2R1W 或 2R2W \\
\bottomrule
\end{longtable}

\hwkey{寄存器堆}（Register File, RF）是 SRAM 的特殊应用：CPU 的乱序执行要求同一拍读出多个源操作数并写回结果，因此寄存器堆采用多端口单元，端口配置用“2R1W”（两读一写）这类记法标注，每多一个端口就要多一对位线和一条字线。以 ARM Cortex-A78 为例，其整数物理寄存器堆的项数与端口组织未见官方公开，第三方微结构分析的通行估计为 128 项量级、2R1W 或 2R2W 配置。由于每个端口都需要独立的位线对和字线，2R1W 寄存器堆的单元面积约为单端口 SRAM 的 \hwevidence{2--3 倍}。

\hwkey{冗余与修复}：

制造缺陷（颗粒、光刻缺陷）会导致某些行或列的单元失效。若整颗宏因此报废，良率损失无法接受，所以 SRAM 宏内嵌\hwkey{冗余行}（spare row）和\hwkey{冗余列}（spare column）：出厂测试发现坏行、坏列后，用激光或电流熔断\term{熔丝}{fuse}（或反其道行之，把\term{反熔丝}{antifuse}击穿成永久导通），将坏地址永久重定向到冗余单元。熔丝与反熔丝都是一次性可编程（one-time programmable, OTP）的物理开关，编程后不可恢复。

\begin{longtable}{L{0.24\linewidth}L{0.26\linewidth}L{0.40\linewidth}}
\toprule
冗余类型 & 典型数量 & 修复方式 \\
\midrule
冗余行（spare row） & 2--4 行 / 阵列 & 熔丝（fuse）或反熔丝（antifuse）将坏行地址重定向到 spare row \\\tablerowrule
冗余列（spare column） & 4--8 列 / 阵列 & 同理；spare column 通常按 byte 或 ECC word 粒度 \\\tablerowrule
ECC 位 & 每 64 数据位 + 8 校验位 & SEC-DED（单错纠正、双错检测）；嵌入阵列中 \\
\bottomrule
\end{longtable}

\hwevidence{典型配置}：一个 64 KiB L1 cache 阵列 = 1024 行 $\times$ 512 列数据 + 64 列 ECC + 4 行冗余 + 8 列冗余。

%======================================================================
\topic{时序参数}
%======================================================================

前面各节讨论的都是阵列内部的工作过程，本小节转向对外接口：SRAM 宏作为可集成的模块交付使用时，使用者必须满足哪些时序约定。SRAM 宏的时序参数定义了外部接口（地址、数据、控制信号）的时序约束：

\begin{longtable}{L{0.22\linewidth}L{0.18\linewidth}L{0.50\linewidth}}
\toprule
参数 & 典型值（28nm） & 定义 \\
\midrule
\hwkey{$t_{AA}$}（地址访问时间） & 0.8--1.5 ns & 地址有效 → 数据输出稳定（读） \\\tablerowrule
\hwkey{$t_{RC}$}（读周期时间） & 1.0--2.0 ns & 两次连续读操作的最小间隔 \\\tablerowrule
\hwkey{$t_{WC}$}（写周期时间） & 1.0--2.0 ns & 写操作的最小周期（地址+数据+WE 有效） \\\tablerowrule
\hwkey{$t_{AS}$}（地址建立时间） & 0.1--0.3 ns & 地址必须在 CLK 上升沿前稳定的最小时间 \\\tablerowrule
\hwkey{$t_{AH}$}（地址保持时间） & 0.05--0.1 ns & CLK 上升沿后地址必须保持的最小时间 \\\tablerowrule
\hwkey{$t_{DS}$}（数据建立时间） & 0.1--0.3 ns & 写数据必须在 CLK 上升沿前稳定 \\\tablerowrule
\hwkey{$t_{DH}$}（数据保持时间） & 0.05--0.1 ns & CLK 上升沿后写数据必须保持 \\\tablerowrule
\hwkey{$t_{WS}$}（WL 建立到 SA 触发） & 0.5--1.0 ns & WL 有效到感放输出的内部延迟 \\
\bottomrule
\end{longtable}

建立时间与保持时间不是人为规定，而是内部物理过程的反映。地址变化后，行译码器的多级门和长字线的 RC 延迟都需要时间才能稳定：若地址在时钟沿前不足 $t_{AS}$ 才变化，译码输出尚未就位，被打开的可能是错误的字线；若时钟沿后地址立刻变化（违反 $t_{AH}$），旧字线尚未关断，可能出现两条字线同时导通的争用。$t_{AA}$ 则是整条读路径——地址缓冲、译码、字线驱动、位线放电、感放、输出级——延迟之和，直接决定 SRAM 宏能工作的最高频率。

\hwevidence{实际数据}：TSMC 28nm HPC+ 工艺的 \hwevidence{512Kb SRAM 宏}，typical corner（0.9V, 25°C）下 $t_{AA}$ = \hwevidence{0.82 ns}（即 1.2 GHz 单周期访问）；worst corner（0.81V, 125°C）下 $t_{AA}$ = \hwevidence{1.45 ns}。

\begin{pdfdiagram}[x=1cm,y=1cm]
  % Timing diagram for read
  \draw[BookRule, line width=0.5pt, ->] (0,0) -- (11,0) node[right, hw label] {时间};
  % CLK
  \draw[BookRule, line width=0.6pt] (0.5,4.0) -- (1.5,4.0) -- (1.5,4.6) -- (3.5,4.6) -- (3.5,4.0) -- (5.5,4.0) -- (5.5,4.6) -- (7.5,4.6) -- (7.5,4.0) -- (9.5,4.0) -- (9.5,4.6) -- (10.5,4.6);
  \node[hw label, anchor=east] at (0.4,4.3) {CLK};
  % ADDR
  \draw[BookAccent, line width=0.6pt] (0.5,2.8) -- (1.0,2.8) -- (1.0,3.4) -- (5.0,3.4) -- (5.0,2.8) -- (5.5,2.8) -- (5.5,3.4) -- (9.5,3.4) -- (9.5,2.8) -- (10.5,2.8);
  \node[hw label, anchor=east] at (0.4,3.1) {ADDR};
  \node[hw label, font=\tiny] at (3.0,3.6) {A0};
  \node[hw label, font=\tiny] at (7.5,3.6) {A1};
  % DOUT
  \draw[BookTeal, line width=0.6pt] (0.5,1.5) -- (2.5,1.5) -- (2.7,1.5) -- (2.7,2.1) -- (6.5,2.1) -- (6.5,1.5) -- (6.7,1.5) -- (6.7,2.1) -- (10.5,2.1);
  \node[hw label, anchor=east] at (0.4,1.8) {DOUT};
  \node[hw label, font=\tiny] at (4.5,2.3) {D(A0)};
  \node[hw label, font=\tiny] at (8.5,2.3) {D(A1)};
  % tAA annotation: 地址有效(ADDR 变为 A0 的沿) → DOUT 输出有效(翻转为 D(A0) 的沿)
  \draw[BookRed, line width=0.5pt] (1.0,0.85) -- (1.0,1.15);
  \draw[BookRed, line width=0.5pt] (2.7,0.85) -- (2.7,1.15);
  \draw[BookRed, line width=0.5pt, <->] (1.0,1.0) -- (2.7,1.0);
  \node[BookRed, font=\tiny\bfseries] at (1.85,0.7) {$t_{AA}$};
  % tAS annotation
  \draw[BookAmber, line width=0.5pt, <->] (1.0,4.9) -- (1.5,4.9);
  \node[BookAmber, font=\tiny\bfseries] at (1.25,5.1) {$t_{AS}$};
  % Caption
  \node[hw label, text width=9cm] at (5.5,-0.5) {SRAM 读时序：地址在 CLK 上升沿前 $t_{AS}$ 建立，数据在 CLK 后 $t_{AA}$ 输出。流水线中 $t_{AA}$ 必须小于一个时钟周期。};
\end{pdfdiagram}
\diagramnote{SRAM 同步读时序。$t_{AA}$ 是限制最高工作频率的关键参数：$f_{max} = 1/t_{AA}$。}

%======================================================================
\topic{功耗：动态与漏电}
%======================================================================

单元结构与阵列组织确定之后，下一个工程问题是功耗——尤其在大容量 cache 占去芯片一半以上面积的今天。SRAM 功耗分为\hwkey{动态功耗}（读写切换）和\hwkey{静态功耗}（漏电）：

\hwkey{动态功耗}：

$$P_{dynamic} = \alpha \cdot C_{BL} \cdot V_{DD}^2 \cdot f + C_{WL} \cdot V_{DD}^2 \cdot f + C_{internal} \cdot V_{DD}^2 \cdot f$$

其中 $\alpha$ 是活动因子（读操作时约 1/2 的 BL 对放电），$f$ 是访问频率。

\begin{longtable}{L{0.28\linewidth}L{0.22\linewidth}L{0.40\linewidth}}
\toprule
功耗来源 & 典型值（28nm, 1GHz） & 说明 \\
\midrule
BL 充放电（读） & 0.5--2 mW/Mbit & 每次读操作 BL 对放电 $\Delta V \times C_{BL}$ \\\tablerowrule
BL 驱动（写） & 1--4 mW/Mbit & 写时 BL 对全摆幅切换 \\\tablerowrule
WL 充放电 & 0.2--0.5 mW/Mbit & 每次访问一条 WL 充放电 \\\tablerowrule
SA + 逻辑 & 0.3--1 mW/Mbit & 感放、译码、mux 切换 \\
\bottomrule
\end{longtable}

\hwkey{静态功耗（漏电）}：

\begin{longtable}{L{0.28\linewidth}L{0.22\linewidth}L{0.40\linewidth}}
\toprule
漏电来源 & 典型值（28nm, 0.9V） & 温度依赖 \\
\midrule
亚阈值漏电（NMOS 截止） & 1--10 pA/管 & $\propto e^{T}$，每升 10°C 约翻倍 \\\tablerowrule
栅极隧穿漏电 & 0.1--1 pA/管 & 弱温度依赖 \\\tablerowrule
PN 结漏电 & 0.5--2 pA/结 & $\propto e^{T}$ \\\tablerowrule
总单元漏电（6 管） & 5--50 pA/cell & — \\\tablerowrule
1 Mbit 阵列总漏电 & 5--50 $\mu$A & 28nm, 0.9V, 25°C \\
\bottomrule
\end{longtable}

\hwevidence{实际数据}：一个 8 MiB L3 cache（64 Mibit，28nm）的静态漏电约 \hwevidence{320 $\mu$A--3.2 mA}（即约 0.3--2.9 mW at 0.9V）。亚阈值漏电本身每升 10°C 约翻倍，但高温下受反型层电荷、体效应等因素影响，总单元漏电从 25°C 到 125°C 典型增大约 \hwevidence{1--2 个数量级}。

\hwkey{功耗管理技术}：

\begin{longtable}{L{0.22\linewidth}L{0.38\linewidth}L{0.30\linewidth}}
\toprule
技术 & 原理 & 效果 \\
\midrule
\hwkey{Power gating} & 用 header/footer 开关切断 VDD 或 VSS & 漏电降为 $\sim$0；但数据丢失 \\\tablerowrule
\hwkey{Retention mode} & 降低 VDD 到 DRV 以上（如 0.5V） & 漏电降低 5--10$\times$；数据保持 \\\tablerowrule
\hwkey{DVFS} & 根据负载动态调节 VDD 和频率 & $P \propto V_{DD}^2 \cdot f$；降压 20\% → 功耗降 36\% \\\tablerowrule
\hwkey{Bitline floating} & 不访问时 BL 浮空（不预充） & 减少 BL 漏电；但下次访问需额外预充时间 \\
\bottomrule
\end{longtable}

表中的 header/footer 开关指串在电源轨与阵列之间的大尺寸开关管：header 管串在 VDD 一侧，footer 管串在 VSS 一侧，关断后阵列与电源物理隔离，漏电只剩开关管自身的截止电流，代价是数据全部丢失，醒来后要重新加载。Retention mode 利用的正是 1.2 节的 DRV 特性：VDD 只要高于 DRV 数据就不丢，而亚阈值漏电随电压近似指数下降，降压的收益极大。\term{动态电压频率调节}{dynamic voltage and frequency scaling, DVFS}则利用 $P\propto V_{DD}^2 f$：频率需求低时连电压一起降，比单独降频多省一份电压平方的收益。

%======================================================================
\topic{软错误：粒子翻转与 ECC}
%======================================================================

前面所有分析都假定环境是安静的；事实上 SRAM 时刻暴露在辐射环境中。本小节讨论高能粒子撞击引起的数据翻转，以及为什么大容量 SRAM 离不开纠错码。\hwkey{软错误}（Soft Error）是由高能粒子（alpha 粒子、宇宙射线中子）击中 SRAM 单元，在存储节点产生足够多的电子-空穴对，导致存储值翻转。这不是器件损坏（硬错误），而是数据错误——重写即可恢复。

\hwkey{物理机制}：
\begin{enumerate}
  \item \hwkey{Alpha 粒子}：即氦原子核，来自封装材料（焊料、引线框架）中的微量 U/Th 衰变，能量 4--9 MeV
  \item \hwkey{宇宙射线中子}：高能中子（$>$1 MeV）与 Si 原子核发生核反应，产生重离子反冲
  \item 粒子在 Si 中产生电子-空穴对（每 3.6 eV 产生一对）
  \item 电荷被存储节点收集；若收集电荷 $>$ \hwkey{临界电荷 $Q_{crit}$}，节点电压被拉过反相器阈值 → 翻转
\end{enumerate}

\hwkey{临界电荷 $Q_{crit}$}：

$$Q_{crit} = C_{node} \times (V_{DD} - V_{th}) \approx C_{cell} \times \Delta V_{SNM}$$

\begin{longtable}{L{0.16\linewidth}L{0.18\linewidth}L{0.18\linewidth}L{0.38\linewidth}}
\toprule
工艺 & $C_{node}$ & $Q_{crit}$ & 软错误率趋势 \\
\midrule
130 nm & $\sim$5 fF & $\sim$3--5 fC & 基准 \\\tablerowrule
65 nm & $\sim$2 fF & $\sim$1.5--3 fC & 约 2$\times$ \\\tablerowrule
28 nm & $\sim$1 fF & $\sim$0.8--1.5 fC & 约 4--5$\times$ \\\tablerowrule
7 nm & $\sim$0.5 fF & $\sim$0.3--0.7 fC & 约 8--10$\times$ \\
\bottomrule
\end{longtable}

\hwkey{错误率度量}——\hwkey{FIT}（Failure In Time）：1 FIT = 每 $10^9$ 器件小时出现 1 次错误。

\begin{longtable}{L{0.24\linewidth}L{0.26\linewidth}L{0.40\linewidth}}
\toprule
条件 & 原始 FIT / Mbit & 说明 \\
\midrule
地面（海平面） & 200--1000 FIT/Mbit & 宇宙射线中子为主 \\\tablerowrule
高海拔（3000m） & 1000--5000 FIT/Mbit & 中子通量增大 5--10$\times$ \\\tablerowrule
航空（10km） & 5000--20000 FIT/Mbit & 屏蔽减少，中子通量极大 \\\tablerowrule
封装 alpha 贡献 & 50--200 FIT/Mbit & 低 alpha 封装材料可降低 \\
\bottomrule
\end{longtable}

关于海拔的影响需要补充定量背景：大气中子通量随海拔近似指数增长，JEDEC JESD89 与 IEC 62396 系列标准给出的量级是 10--12 km 巡航高度的中子通量约为海平面的 100--300 倍，因此航空电子的软错误问题远比地面严峻。表中 FIT 数值是不同文献来源综合出的量级参考，实际错误率还随工艺、电压、封装材料和纬度显著变化，不宜当作精确指标使用。

\hwevidence{实际影响}：一个 128 MiB L3 cache（1024 Mibit）在海平面的原始软错误率约 \hwevidence{200,000--1,000,000 FIT}（即平均每 \hwevidence{1--7 个月}出现一次未纠正错误）。必须使用 ECC。

\hwkey{ECC 保护}——\hwkey{SEC-DED}（Single Error Correction, Double Error Detection）：

\begin{longtable}{L{0.24\linewidth}L{0.26\linewidth}L{0.40\linewidth}}
\toprule
ECC 方案 & 开销 & 能力 \\
\midrule
奇偶校验（parity） & +1 bit / 8 data & 仅检测单 bit 错误，不能纠正 \\\tablerowrule
SEC-DED (Hamming) & +8 bit / 64 data & 纠正 1 bit，检测 2 bit \\\tablerowrule
SEC-DED + 符号纠错 & +16 bit / 128 data & 纠正 1 bit + 部分多 bit \\\tablerowrule
Chipkill / SDDC & +16 bit / 128 data & 纠正整颗芯片失效（服务器） \\
\bottomrule
\end{longtable}

SEC-DED 的经典实现是 Hamming 码：编码时把每个数据位编入若干奇偶校验组，读出时重算各组奇偶并拼成\term{校正子}{syndrome}。校正子为零表示无错；非零时其数值直接指出出错位的位置，翻转该位即完成纠正。再加一位覆盖全字的全局奇偶校验，就能区分“一位错”（校正子非零且全局校验不符，可纠）与“两位错”（校正子非零但全局校验相符，只能报警），这就是 DED 能力的来源。

\hwkey{ECC 在 SRAM 中的实现}：每 64 数据位附加 8 校验位（共 72 bit），校验位与数据位一起存储在阵列中。读操作时，72 bit 一起读出，ECC 逻辑在感放输出后 1--2 个门延迟内完成校验和纠正。写操作时，ECC 编码器先计算校验位，再一起写入。

\begin{lstlisting}
// 注释（推进）：每轮只推进一个元素或一个控制步，并保持中间状态的一致性。
// 注释（边界）：退出条件成立后才提交结果；空集合、越界和失败要单独处理。
ECC SEC-DED for 64-bit data word:
  // Encode (write path):
  parity[0] = XOR(data[0], data[1], data[3], data[4], data[6], ...)
  parity[1] = XOR(data[0], data[2], data[3], data[5], data[6], ...)
  ...
  parity[6] = XOR(data[0..63] selected by Hamming matrix)
  parity[7] = XOR(data[0..63], parity[0..6])  // overall parity for DED

  // Decode (read path):
  syndrome[6:0] = H_matrix * received_word[71:0]
  if syndrome == 0: no error
  if syndrome != 0 and parity[7] mismatch: single-bit error -> flip bit[syndrome]
  if syndrome != 0 and parity[7] match: double-bit error -> flag uncorrectable
\end{lstlisting}

%======================================================================
\topic{真实芯片中的 SRAM 参数}
%======================================================================

前面各节的数值多来自教材与工艺文献的典型范围，本小节给出真实产品的参数，与前面的典型范围相互印证。

\hwkey{Intel 处理器 Cache 参数}（公开数据）：

\begin{longtable}{L{0.19\linewidth}L{0.12\linewidth}L{0.12\linewidth}L{0.12\linewidth}L{0.12\linewidth}L{0.22\linewidth}}
\toprule
处理器 & L1D & L1I & L2 & L3 & 工艺 \\
\midrule
Intel Core i9-13900K & 48 KiB/核 & 32 KiB/核 & 2 MiB/核 & 36 MiB & Intel 7 (10nm) \\\tablerowrule
Intel Core i9-14900K & 48 KiB/核 & 32 KiB/核 & 2 MiB/核 & 36 MiB & Intel 7 \\\tablerowrule
Intel Xeon w9-3495X & 48 KiB/核 & 32 KiB/核 & 2 MiB/核 & 105 MiB & Intel 7 \\\tablerowrule
Intel Core Ultra 200S & 48 KiB/核 & 32 KiB/核 & 3 MiB/核 & 36 MiB & Intel 20A/TSMC N3 \\
\bottomrule
\end{longtable}

表中 Core Ultra 200S（Arrow Lake）一栏的工艺标注需要说明：这一代处理器的计算模块（compute tile）实际由台积电 N3B 代工，Intel 20A 用于同期其他产品线；表中“Intel 20A/TSMC N3”反映的是该代产品处于自家工艺与外包工艺并用的过渡期，阅读时按“3nm 级”理解即可。

\hwkey{ARM Cortex 系列寄存器堆与 Cache}：

\begin{longtable}{L{0.22\linewidth}L{0.18\linewidth}L{0.18\linewidth}L{0.32\linewidth}}
\toprule
核 & 整数 RF & L1D & 备注 \\
\midrule
Cortex-A78 & 128 $\times$ 64 bit (2R1W) & 64 KiB & 13 级流水线 \\\tablerowrule
Cortex-A710 & 128 $\times$ 64 bit (2R1W) & 64 KiB & Armv9 \\\tablerowrule
Cortex-X3 & 192 $\times$ 64 bit (2R2W) & 64 KiB & 大核，更宽发射 \\\tablerowrule
Cortex-X4 & 192+ $\times$ 64 bit & 64 KiB & Armv9.2 \\\tablerowrule
Neoverse V2 & 192 $\times$ 64 bit & 64 KiB & 服务器核 \\
\bottomrule
\end{longtable}

需要说明数据来源：ARM 不公开寄存器堆的物理项数与端口组织，表中整数 RF 一栏综合自第三方微结构分析的实测与推断，属于业界通行估计而非官方规格；L1D 容量与流水线级数为公开数据。

\hwkey{SRAM 宏参数（来自工艺设计套件 PDK（process design kit）与 Foundry 公开数据）}：

\begin{longtable}{L{0.18\linewidth}L{0.14\linewidth}L{0.14\linewidth}L{0.14\linewidth}L{0.30\linewidth}}
\toprule
参数 & TSMC 28nm & TSMC 7nm & Samsung 5nm & 说明 \\
\midrule
最小单元面积 & 0.126 $\mu$m$^2$ & 0.027 $\mu$m$^2$ & 0.021 $\mu$m$^2$ & 6T 单元（含接触） \\\tablerowrule
宏容量范围 & 4Kb--32Mb & 4Kb--256Mb & 4Kb--256Mb & 可配置 \\\tablerowrule
最高频率 (typ) & 1.5--2.0 GHz & 2.5--3.5 GHz & 3.0--4.0 GHz & 取决于容量和配置 \\\tablerowrule
读电流 (active) & 0.5--2 mA/Mb & 0.3--1.5 mA/Mb & 0.2--1.2 mA/Mb & VDD=nominal \\\tablerowrule
漏电 (standby) & 10--50 $\mu$A/Mb & 20--100 $\mu$A/Mb & 30--150 $\mu$A/Mb & 25°C \\\tablerowrule
冗余 & 2--4 spare row/col & 4--8 spare row/col & 4--8 spare row/col & 修复良率 \\\tablerowrule
Assist 电路 & 可选 NBL & NBL + WL underdrive & VDD collapse + BB & 先进节点标配 \\
\bottomrule
\end{longtable}

需要说明：表中 5nm 列的单元面积 0.021 $\mu$m$^2$ 是公开发表的 TSMC N5 高密度单元数据（Samsung 5LPE 未见同口径的公开数值，报道量级相近）；该列其余参数取 5nm 级工艺的代表性范围，用于展示代际趋势，不宜当作某一家的精确规格引用。

\hwevidence{TSMC 7nm SRAM 密度}：\hwevidence{0.027 $\mu$m$^2$/bit}，即 1 Mbit 阵列面积约 \hwevidence{27,000 $\mu$m$^2$}（约 0.027 mm$^2$）。一个 64 MiB（512 Mbit）L3 cache 的纯阵列面积约 \hwevidence{13.8 mm$^2$}，加上外围电路（SA、译码器、驱动器）总面积约 \hwevidence{20--25 mm$^2$}。

\hwkey{SRAM 面积不缩放的困境}：

\begin{longtable}{L{0.16\linewidth}L{0.20\linewidth}L{0.20\linewidth}L{0.34\linewidth}}
\toprule
节点 & 单元面积 & 缩放比 & 问题 \\
\midrule
28 nm & 0.126 $\mu$m$^2$ & 1.0$\times$ & 基准 \\\tablerowrule
14 nm & 0.064 $\mu$m$^2$ & 0.51$\times$ & 接近理想（0.5$\times$） \\\tablerowrule
7 nm & 0.027 $\mu$m$^2$ & 0.42$\times$ & 开始偏离理想 \\\tablerowrule
5 nm & 0.021 $\mu$m$^2$ & 0.78$\times$ & 严重不缩放（fin 数限制） \\\tablerowrule
3 nm & 0.019 $\mu$m$^2$ & 0.90$\times$ & 几乎不缩放 \\
\bottomrule
\end{longtable}

从 7nm 到 5nm 到 3nm，SRAM 单元面积缩放比从 0.42$\times$ 退化到 0.90$\times$（几乎不缩小）。这是因为 FinFET 的 fin 数只能取整数（最小 1 fin），而 6T 单元需要至少 6 个 fin（每管 1 fin），加上接触和连线规则，单元宽度无法继续缩小。\hwevidence{这是先进节点 die 面积中 SRAM 占比越来越大的根本原因}——以 Apple M2 为例，第三方 die 照片分析显示其系统级缓存（system-level cache, SLC，公开报道容量约 24 MiB）连同外围电路占去了 die 面积中相当可观的比例（估计在 10\% 以上）。

\begin{warningbox}
\hwkey{常见误区}：很多人认为“SRAM 就是快一点的 DRAM”。实际上两者的物理机制完全不同：SRAM 靠正反馈锁存（6 管），只要通电就保持；DRAM 靠电容存储（1 管 1 电容），必须每 64 ms 刷新。SRAM 的读是非破坏性的，DRAM 的读是破坏性的。SRAM 的失效模式是 SNM 不足导致软错误，DRAM 的失效模式是电容漏电导致数据丢失。两者不能互换。
\end{warningbox}

\section{二、Cache 的 tag、index 与 offset}

Cache 用地址低位选择行内字节，用中间位选择集合，用高位 tag 判断当前集合里保存的是哪段主存。以 32 KiB、64 B cache line、8 路组相联、48-bit 物理地址为例：共有 $32\,\text{KiB}/(64\times8)=64$ 个集合，offset 为 6 bit，index 为 6 bit，其余 36 bit 是 tag。

一次读取同时访问所选集合的 8 份 tag 和数据阵列。tag 比较命中且 valid 为 1 时，由路选择器取出对应字节；全部不命中时，Cache 控制器分配 miss 状态项，向下一级请求整条 64 B line。若牺牲行 dirty，还要先写回。

\begin{longtable}{L{0.2\linewidth}L{0.32\linewidth}L{0.38\linewidth}}
\toprule
策略 & 硬件动作 & 必须明确的边界 \\
\midrule
write-through & 每次写同时送往下一级 & 写缓冲何时满，完成代表哪一级完成 \\\tablerowrule
write-back & 只改本级并置 dirty & 替换、同步或持久化时何时写回 \\\tablerowrule
write-allocate & 写 miss 先取整行再修改 & 额外读流量与所有权获取 \\\tablerowrule
no-write-\allowbreak allocate & 写 miss 直接送下一级 & 部分写合并与字节使能 \\
\bottomrule
\end{longtable}

\section{三、多级 Cache、一致性与内存序}

L1 追求每拍访问，容量小且常分成指令与数据两套；L2 容量更大；末级 Cache 还承担核心间共享与外部内存流量过滤。平均访问时间可写成

\[
T_{avg}=T_{L1}+m_{L1}(T_{L2}+m_{L2}(T_{LLC}+m_{LLC}T_{DRAM})).
\]

多核中，同一物理 cache line 可能被多个核心保存。写者必须先取得独占所有权，使其他副本失效；读者若命中别处的修改副本，数据要由拥有者或已更新的下一级返回。一致性解决“同一地址最终看到同一写入次序”，内存序还要规定不同地址之间哪些重排对软件可见，二者不能混为一谈。

DMA 设备若不参与硬件一致性，驱动必须在所有权交接处清理或失效 Cache，并使用屏障保证“描述符内容先可见，doorbell 后可见”。这条规则会在设备事务章节再次出现。

\section{四、命中路径：一拍之内完成哪些工作}

处理器给出虚拟地址后，L1 Cache 不能只做一次“查表”。在典型的虚拟索引、物理标记设计中，页内偏移里的 index 先选择集合；TLB 同时给出物理页号；随后将物理 tag 与各路 tag、valid 和权限状态比较。命中的一路还要经过字节选择、对齐检查和转发选择，最终才能把数据送到执行单元。若 load/store queue 中存在更早的同地址写，数据甚至可能来自写转发，而不是 SRAM 数据阵列。

这条并行路径解释了为什么 L1 容量、路数和时钟频率不能独立选择。增加路数能够减少冲突 miss，却会增加 tag 比较器数量和数据多路选择器延迟；增大 line 会减少 tag 数量，却会带来更大的误取流量。工程上常把 tag 阵列、数据阵列和路预测结合起来，使常见命中保持短延迟，并为预测错误保留重选路径。

\begin{pdfdiagram}[
  node distance=8mm and 9mm,
  cachebox/.style={draw, rounded corners, align=center, minimum width=25mm, minimum height=9mm, fill=blue!5},
  cachesmall/.style={draw, rounded corners, align=center, minimum width=20mm, minimum height=8mm, fill=green!6},
  >=Stealth]
\node[cachebox] (va) {虚拟地址};
\node[cachesmall, right=of va, yshift=8mm] (idx) {index\\选择集合};
\node[cachesmall, right=of va, yshift=-8mm] (tlb) {TLB\\产生物理 tag};
\node[cachebox, right=of idx] (arrays) {多路数据阵列\\并行读出};
\node[cachebox, right=of tlb] (tags) {多路 tag 比较\\valid/权限检查};
\node[cachebox, right=of arrays, yshift=-8mm] (mux) {命中路选择\\字节与对齐选择};
\node[cachebox, right=of mux] (out) {数据返回或\\异常/缺失};
\draw[->] (va) -- (idx);
\draw[->] (va) -- (tlb);
\draw[->] (idx) -- (arrays);
\draw[->] (tlb) -- (tags);
\draw[->] (arrays) -- (mux);
\draw[->] (tags) -- (mux);
\draw[->] (mux) -- (out);
\end{pdfdiagram}
\diagramnote{一次 L1 读取的并行关键路径。集合选择与 TLB 翻译同时开始，物理 tag 比较结果再选择对应数据路；图中上下分支之间保留间距，以区分数据阵列和翻译比较两条路径。}

\section{五、缺失路径：MSHR、写回缓冲与背压}

阻塞式 Cache 在一次 miss 完成前停止接受新的访问，结构简单但会浪费主存等待时间。非阻塞 Cache 用 MSHR（miss status holding register）记录每个未完成 line 的地址、请求者、访问大小、目标寄存器和权限状态。后续请求若命中同一个在途 line，可以合并到原 MSHR；命中其他常驻 line 的访问则继续执行，这分别称为“hit under miss”和“miss under miss”。

一次替换通常包含以下状态，而不是一个瞬时动作：选择 victim；若 victim 为 dirty，将其地址和数据放入写回缓冲；为新 line 分配 MSHR；向下一级发请求；接收若干 beat；在满足关键字优先或整行到齐条件后唤醒等待者；最后提交 tag、valid 和一致性状态。写回缓冲、MSHR 或下行信用任一耗尽，都会沿请求路径产生背压。因此性能计数器中的“Cache miss”还应区分合并 miss、MSHR 满、写回缓冲满和下一级阻塞。

\begin{longtable}{L{0.2\linewidth}L{0.31\linewidth}L{0.39\linewidth}}
\toprule
资源 & 保存的关键状态 & 资源满时的可见结果 \\
\midrule
MSHR & line 地址、等待者、返回 beat、错误状态 & 新 miss 暂停，已有命中仍可能前进 \\\tablerowrule
写回缓冲 & victim 地址、dirty 数据、字节掩码 & 不能继续替换 dirty line \\\tablerowrule
填充缓冲 & 从下一级返回但尚未提交的数据 & 返回通道施加背压或延后新的读取 \\\tablerowrule
store buffer & 尚未写入 Cache 的写及其顺序信息 & 提交端停止，屏障等待时间上升 \\
\bottomrule
\end{longtable}

\section{六、替换、预取与包含关系}

真正的 LRU 需要记录完整的最近使用次序，路数增加后代价迅速上升，因此硬件常用 tree-PLRU、NRU 或随机替换近似。替换策略不能只看命中率：脏行的写回成本、共享行的失效成本、预取行的置信度也会影响 victim 选择。流式访问还可能逐出循环反复使用的工作集，所以“更积极预取”并不总能提高性能。

预取器根据顺序流、固定步长或相关地址提前请求 line。评价它至少要同时看准确率、覆盖率、及时性和污染：有用但太晚的预取仍落在需求 miss 之后；准确却占满 MSHR 的预取会延迟需求请求；错误预取既消耗互连带宽又挤出有效数据。常见控制办法是降低预取优先级、动态调节距离，并在压力升高时暂停发射。

多级 Cache 还要声明包含关系。inclusive 设计让上级 line 必然在下级有副本，便于用末级目录筛选探测，但下级替换可能迫使上级失效；exclusive 设计尽量让各级保存不同 line，提高总有效容量，却增加迁移复杂度；non-inclusive/non-exclusive 不承诺两者，需用独立目录维护共享者信息。包含关系是数据放置策略，不等于一致性协议本身。

\section{七、从 MESI 状态到真实写入事务}

以 MESI 为例，I 表示本地无有效副本，S 表示可能有多个只读副本，E 表示只有本核心持有且数据干净，M 表示本核心持有已修改数据。读取 miss 通常发出共享读；若没有其他共享者，可进入 E，否则进入 S。对 S 状态 line 写入时，核心必须请求升级并等待其他共享者确认失效；只有取得写权限后才能让该 store 对一致性域生效。E 状态的本地写可静默转为 M。

两个核心写同一 line 中不同变量仍会反复争夺所有权，这就是 false sharing。变量在源代码中互不相关并不能阻止 line 粒度的失效。定位时应结合失效探测、远端命中和 line 迁移计数，再通过数据布局或每核分片减少共享写，而不是简单增大 Cache。

一致性事务还要处理竞争和暂态。例如某核心发出升级请求、尚未收到全部失效确认时，line 既不能被当作稳定 S，也不能被当作稳定 M；控制器需保存暂态、待收确认数和可能到达的数据响应。协议死锁分析必须覆盖请求网、响应网和 snoop 网之间的资源依赖。

\section{八、Cache 中的校验与故障闭合}

tag 和状态位常用奇偶校验或 ECC，数据阵列可按若干字节分组保护。可纠正错误在读取时被修正并记录，后台 scrub 再把正确值写回；不可纠正错误不能继续作为普通数据交付。若 dirty line 出现不可纠正错误，下一级没有最新副本，处理器通常需要上报机器检查并隔离受影响的执行上下文。若 clean line 损坏，则可将其失效并从下一级重新获取，但仍应保留诊断信息。

错误处理与正常 miss 共用许多资源，设计时要确保错误响应不会因 MSHR 或写回队列已满而永久等待。可观测性至少包括错误物理地址、阵列与路、可纠正/不可纠正分类、首次错误状态和累计计数；软件据此区分偶发位翻转、持续坏单元与更大故障域。




























\section{九、用 AMAT 把命中率变成实际停顿}

Cache 的价值不能只用“命中率很高”判断，因为不同层 miss 的代价不同。平均访存时间 AMAT（average memory access time）可按条件概率逐层展开。设 L1 命中时间 1 ns，L1 miss 率 5\%；L2 在收到 L1 miss 后的命中时间为 6 ns，局部 miss 率 20\%；主存额外代价为 80 ns，则：
\[
\mathrm{AMAT}=1+0.05\times(6+0.20\times80)=2.1\ \mathrm{ns}.
\]

这里的 20\% 是“到达 L2 的请求中有多少仍 miss”，不是全部访问的全局比例。真正到达主存的比例是 $0.05\times0.20=1\%$。若把 20\% 直接乘到全部访问，会把主存流量高估二十倍。

平均值还不能替代并行性。阻塞式核心可能把一次 80 ns miss 全部计入停顿；乱序核心若有足够 MSHR 和独立工作，可以让多个 miss 重叠。此时 AMAT 仍描述单次访问的平均层次代价，但程序时间还取决于内存级并行度、依赖链和 ROB 能覆盖的窗口。评价 Cache 应同时报告命中率、每级局部 miss 率、miss 延迟、在途数和最终停顿拍数。

\begin{longtable}{L{0.20\linewidth}L{0.25\linewidth}L{0.25\linewidth}L{0.20\linewidth}}
\toprule
指标 & 分母 & 回答的问题 & 常见误用 \\
\midrule
L1 miss 率 & 全部 L1 访问 & 多少请求进入 L2 & 与 L2 局部 miss 率直接相加 \\\tablerowrule
L2 局部 miss 率 & 到达 L2 的请求 & L1 miss 中多少进入主存 & 当成全部访问比例 \\\tablerowrule
MPKI & 每千条退休指令 & 程序产生多少 miss & 忽略每个 miss 的代价和重叠 \\\tablerowrule
停顿周期 & 总周期或每条指令 & miss 最终阻塞了多少执行 & 只由 miss 数推算而不看并行性 \\
\bottomrule
\end{longtable}

\topic{逐地址走完一次冲突 miss}

设直接映射 Cache 总容量 1 KiB，line 为 64 B，共 16 个集合，index 占 4 bit，offset 占 6 bit。地址 \code{0x0000} 与 \code{0x0400} 相差 1024 B，offset 和 index 完全相同，只是 tag 不同，因此二者都映射到集合 0。

\begin{longtable}{L{0.11\linewidth}L{0.17\linewidth}L{0.16\linewidth}L{0.16\linewidth}L{0.30\linewidth}}
\toprule
次序 & 地址 & 集合 & tag & 结果与集合 0 状态 \\
\midrule
1 & \code{0x0000} & 0 & 0 & miss，装入 tag 0 \\\tablerowrule
2 & \code{0x0400} & 0 & 1 & miss，替换 tag 0，装入 tag 1 \\\tablerowrule
3 & \code{0x0000} & 0 & 0 & miss，替换 tag 1 \\\tablerowrule
4 & \code{0x0400} & 0 & 1 & miss，再次替换 tag 0 \\
\bottomrule
\end{longtable}

工作集只有两条 line，远小于 1 KiB 容量，仍然每次 miss；这不是容量 miss，而是映射冲突。改成 2 路组相联并保持总容量不变后，集合数降为 8，两个地址仍进入同一集合，但可以分别占据两路；前两次强制 miss 后，后两次变为 hit。代价是每次访问要并行读两路 tag 和数据，再由比较结果选择一路。

这个例子也给出定位方法：先把地址按 line 大小右移得到 line 编号，再对集合数取模得到 index，最后比较同时活跃地址是否反复落到少数集合。只看总工作集大小无法发现冲突；性能计数器若支持按集合或物理地址采样，可以进一步确认是容量不足、冲突还是共享失效。

\topic{非阻塞 Cache 的暂态与返回 beat}

稳定的 MESI 字母不足以描述 miss 正在进行时的状态。某 line 从 I 发出共享读、等待数据时，可以记为 IS；从 S 发出升级、等待其他核心失效确认时，可以记为 SM；M line 被替换、等待写回完成时也需要独立暂态。暂态项至少保存目标稳定状态、待收响应数量、已经返回的 beat、等待的处理器请求和错误标记。

\begin{longtable}{L{0.13\linewidth}L{0.22\linewidth}L{0.28\linewidth}L{0.27\linewidth}}
\toprule
暂态 & 触发动作 & 等待的事件 & 完成或失败边界 \\
\midrule
IS & I 状态 load miss 发共享读 & 数据响应与权限结果 & 数据到齐后进入 S/E；错误则唤醒异常 \\\tablerowrule
IM & I 状态 store miss 请求独占 & 数据与全部失效确认 & 两者都满足后进入 M \\\tablerowrule
SM & S 状态写请求升级 & 其他共享者确认失效 & 确认数归零后允许写生效 \\\tablerowrule
MI & dirty victim 被替换 & 下级接受完整写回 & 写回完成后释放 victim 资源 \\
\bottomrule
\end{longtable}

若一条 64 B line 经 16 B 总线分四个 beat 返回，Cache 可以采用 critical-word-first：包含处理器所需字的 beat 到达后先唤醒 load，其他 beat 继续填充。但 tag 和 line 状态在整行未闭合前仍处于暂态；另一个访问不能把“关键字已返回”误认为“整行已稳定”。若后续 beat 报错，控制器还要按体系结构规定撤销尚未提交的使用者或报告机器检查。

MSHR 合并也要比较访问属性。两个普通 load 命中同一在途 line 通常可以合并；一个取独占的 store 到达已有共享读 miss 时，可能升级原事务，也可能等待读完成后再发升级，取决于协议。对 MMIO、不可缓存访问或不同权限域的请求，不能只因地址相同就合并。

\section{十、目录一致性把共享者集合变成显式状态}

小规模共享总线可以广播 snoop：每个核心都观察请求并检查本地 tag。核心数增加后，广播带宽和比较功耗上升，目录协议改为在 home 节点记录某条 line 的所有者和共享者集合。读 miss 先到 home；若 line 为 clean shared，home 返回数据并加入请求者；若某核心持有 M，home 转发请求给 owner，由 owner 提供最新数据并按请求类型降级或交出所有权。

以两个核心 C0、C1 和 home H 为例，初始 line A 在 C0 为 M：

\begin{longtable}{L{0.10\linewidth}L{0.23\linewidth}L{0.29\linewidth}L{0.28\linewidth}}
\toprule
步骤 & 消息 & 状态变化 & 为什么不能提前完成 \\
\midrule
1 & C1 向 H 发 ReadShared(A) & H 查到 owner=C0 & H 的内存副本可能是旧值 \\\tablerowrule
2 & H 向 C0 发 ForwardRead(A) & C0 准备最新数据 & C1 尚未收到数据和权限 \\\tablerowrule
3 & C0 向 C1/H 返回数据 & C0 从 M 降到 S，H 更新目录 & 目录更新与数据响应要对应同一事务 \\\tablerowrule
4 & C1 收到数据并确认 & C1 进入 S & 此时 load 才能以正确数据完成 \\
\bottomrule
\end{longtable}

若 C1 请求写权限，H 还要向所有共享者发送 invalidate 并等待确认。目录项的共享者表示可以用完整位图、有限指针或分层目录；位图随核心数增长，有限指针溢出时可能退化为广播。无论编码方式如何，协议必须保证任一时刻至多一个可写所有者，并保证最新 dirty 数据不会因 owner 失联或目录错误而丢失。

消息网络也会死锁。请求、转发、数据和确认若共享有限缓冲，可能形成“请求占满队列等待响应，响应又没有队列空间”的环。实现通常使用不同虚拟通道或规定严格的资源获取次序，并保证某些响应通道始终可前进。一致性正确性不仅是 MESI 状态图，还包括消息、缓冲和互连的前进性。

\topic{伪共享的量化代价}

设两个核心各自每 20 ns 更新一个 8-byte 计数器，但两个计数器位于同一 64 B line。每次写都需要把 line 的 M 权限从另一核心迁移过来；若一次往返所有权迁移为 80 ns，那么单次更新的最低间隔不再由本地 ALU 决定，而接近 80 ns，理论更新率从每核 5000 万次/秒降到全局约 1250 万次/秒量级。

把两个计数器按 64 B 对齐分开放置后，每个核心长期保持自己 line 的 M 状态，更新走本地 L1；最终汇总时才发生共享读取。数据字节数几乎没有变化，性能差异来自一致性权限事务。定位伪共享时应寻找“源代码对象不同、物理 line 相同、写失效和远端 line 命中频繁”这组证据。

\topic{指令 Cache 与自修改代码的同步边界}

数据 Cache 中写入的新机器码不保证立即被指令 Cache 和取指前端观察。某些体系结构维持 I/D Cache 硬件一致，另一些要求软件 clean 数据 Cache、invalidate 指令 Cache，并执行指令同步屏障。动态链接器修补跳转表、JIT 编译器生成代码和调试器写断点都必须遵守平台规定。

正确协议的目标是让三个状态对齐：数据写已到达指令侧可见的一致性点；旧指令副本已经失效；流水线中此前预取或译码的旧字节已经被丢弃。只清 D-Cache 不会删除 I-Cache 旧副本，只清 I-Cache 也不能保证脏数据已经向指令侧传播，只做普通内存屏障更不能替代专用的指令同步边界。

\section*{章末练习}

\begin{chapterexercise}{综合练习：1 至 2}
\begin{exercisesubproblem}{（a）1}
地址分解
\exercisecheck{题目}{32 KiB、64 B line、8 路组相联 Cache 使用 48-bit 物理地址。求集合数以及 offset、index、tag 位数。}
\end{exercisesubproblem}
\begin{exercisesubproblem}{（b）2}
路径分析
\exercisecheck{题目}{按时间顺序说明一次 L1 读取命中的 tag、TLB、数据阵列和路选择动作，并指出增加路数影响哪段关键路径。}
\end{exercisesubproblem}
\end{chapterexercise}

\begin{chapterexercise}{综合练习：3 至 4}
\begin{exercisesubproblem}{（a）3}
状态追踪
\exercisecheck{题目}{同一 line 上先发生两个 load miss，随后另一地址发生 miss。说明 MSHR 合并条件，以及资源不足时的背压位置。}
\end{exercisesubproblem}
\begin{exercisesubproblem}{（b）4}
一致性
\exercisecheck{题目}{两个核心交替修改同一 line 中的两个不同计数器。解释 line 状态如何迁移，并给出一种数据布局改进。}
\end{exercisesubproblem}
\end{chapterexercise}

\begin{chapterexercise}{综合练习：5 至 6}
\begin{exercisesubproblem}{（a）5}
预取评价
\exercisecheck{题目}{某预取器准确率很高，但程序反而变慢。列出至少三项还需检查的指标或资源冲突。}
\end{exercisesubproblem}
\begin{exercisesubproblem}{（b）6}
故障处理
\exercisecheck{题目}{clean line 与 dirty line 分别发生不可纠正错误时，为何恢复策略不同？应记录哪些诊断信息？}
\end{exercisesubproblem}
\end{chapterexercise}

\begin{chapterexercise}{综合练习：AMAT 演算至冲突 trace}
\begin{exercisesubproblem}{（a）AMAT 演算}
L1 命中 1 ns、miss 率 4\%；L2 局部命中时间 8 ns、局部 miss 率 10\%；主存额外代价 90 ns。计算 AMAT，以及全部访问中真正到达主存的比例。
\exercisecheck{学习目标}{逐层使用条件 miss 率}
\end{exercisesubproblem}
\begin{exercisesubproblem}{（b）冲突 trace}
直接映射 Cache 容量 2 KiB、line 64 B。判断 \code{0x0000}、\code{0x0800}、\code{0x1000} 是否映射到同一集合，并走完顺序 A、B、C、A 的命中结果。
\exercisecheck{学习目标}{由 line 编号和集合数识别冲突 miss}
\end{exercisesubproblem}
\end{chapterexercise}

\begin{chapterexercise}{综合练习：暂态闭合至目录事务}
\begin{exercisesubproblem}{（a）暂态闭合}
一条 store miss 已收到数据，但还差两个失效确认。此时能否把 line 标成稳定 M 并让 store 对外可见？还必须保存哪些状态？
\exercisecheck{学习目标}{数据返回和写权限确认是两个完成条件}
\end{exercisesubproblem}
\begin{exercisesubproblem}{（b）目录事务}
某 line 在 C0 为 M，C1 发起读。写出 home、C0、C1 之间的消息和最终状态；说明为什么 home 不能直接从内存返回数据。
\exercisecheck{学习目标}{dirty owner 持有最新副本}
\end{exercisesubproblem}
\end{chapterexercise}

\begin{chapterexercise}{综合练习：伪共享预算至自修改代码}
\begin{exercisesubproblem}{（a）伪共享预算}
两个线程各以 25 ns 周期写同一 line 中不同变量，所有权迁移往返为 100 ns。估算共享 line 能支持的更新率上限，并说明按 line 分离后的主要变化。
\exercisecheck{学习目标}{把一致性迁移时延换算为吞吐上限}
\end{exercisesubproblem}
\begin{exercisesubproblem}{（b）自修改代码}
JIT 把新指令写入普通可写页面后准备跳转执行。列出数据侧、指令侧和流水线侧需要闭合的三个边界。
\exercisecheck{学习目标}{区分数据可见、指令副本失效和前端同步}
\end{exercisesubproblem}
\end{chapterexercise}

\begin{chapterexercise}{综合练习：SRAM 读写周期至 cache 行}
\begin{exercisesubproblem}{（a）SRAM 读写周期}
画出地址、CS、OE、WE、数据线和采样窗口
\exercisecheck{预期结果}{能指出谁驱动数据线，谁采样}
\end{exercisesubproblem}
\begin{exercisesubproblem}{（b）cache 行}
用 tag/index/offset 解释一次 cache 命中和一次 miss
\exercisecheck{预期结果}{能说明 cache line、valid、dirty 和替换边界}
\end{exercisesubproblem}
\end{chapterexercise}

\begin{chapterexercise}{综合练习：组相联查找至替换策略}
\begin{exercisesubproblem}{（a）组相联查找}
用本章的四次访问序列，分别在直接映射和 2 路组相联下走一遍
\exercisecheck{预期结果}{能指出哪几次由 miss 变 hit 及原因}
\end{exercisesubproblem}
\begin{exercisesubproblem}{（b）替换策略}
2 路组相联、LRU，某组初始为空，依次访问 A、B、A、C、A、B；逐步写出命中/失效和组内容，并与同 index 直接映射对比
\exercisecheck{预期结果}{2 路共 3 次 miss，直接映射 6 次全 miss}
\end{exercisesubproblem}
\end{chapterexercise}

\begin{chapterexercise}{预取演算}
顺序读 16 条 cache line，每次 miss 延迟 30 拍、每行处理 8 拍；计算无预取与 next-line 预取的总拍数，并求能完全隐藏预取的最小处理节拍
\exercisecheck{预期结果}{608 拍对 488 拍，处理节拍不小于 30 才能完全隐藏}
\end{chapterexercise}

\begin{chapterexercise}{MESI 走查}
两核初始均为 I，按“A 读 X、B 读 X、A 写 X、B 读 X”写出每步两核状态和互连消息
\exercisecheck{预期结果}{每步最新值位置唯一，失效与写回/转发正确}
\end{chapterexercise}

\section*{参考解答与要点}

\begin{chapteranswer}{综合练习：1 至 2}
\begin{answersubproblem}{（a）1}
地址分解
\answercheck{解答要点}{集合数为 $32\,\text{KiB}/(64\times8)=64$；offset 6 bit，index 6 bit，tag 为 $48-6-6=36$ bit。}
\end{answersubproblem}
\begin{answersubproblem}{（b）2}
路径分析
\answercheck{解答要点}{index 先选集合，TLB 并行给出物理 tag；多路 tag 比较与数据读出并行，命中结果控制路选择和字节选择。路数增加主要加重 tag 比较和数据多选延迟。}
\end{answersubproblem}
\end{chapteranswer}

\begin{chapteranswer}{综合练习：3 至 4}
\begin{answersubproblem}{（a）3}
状态追踪
\answercheck{解答要点}{line 地址相同且请求属性兼容时合并到同一 MSHR；不同 line 需要新项。MSHR 满会阻止新 miss，填充或下行资源不足还会使返回或请求通道背压。}
\end{answersubproblem}
\begin{answersubproblem}{（b）4}
一致性
\answercheck{解答要点}{每次写都要使对方副本失效并取得 M/E 权限，line 在核心间反复迁移。可把计数器放到不同 line，或改为每核计数后汇总。}
\end{answersubproblem}
\end{chapteranswer}

\begin{chapteranswer}{综合练习：5 至 6}
\begin{answersubproblem}{（a）5}
预取评价
\answercheck{解答要点}{检查覆盖率、及时性、Cache 污染、互连带宽、MSHR/队列占用及需求请求优先级；高准确率不保证请求足够及时或资源代价可接受。}
\end{answersubproblem}
\begin{answersubproblem}{（b）6}
故障处理
\answercheck{解答要点}{clean line 可失效后重取；dirty line 的最新值只在本级，通常需上报不可恢复错误。记录地址、路与阵列、错误等级、综合征、首次错误和累计计数。}
\end{answersubproblem}
\end{chapteranswer}

\begin{chapteranswer}{综合练习：AMAT 演算至冲突 trace}
\begin{answersubproblem}{（a）AMAT 演算}
\[
\mathrm{AMAT}=1+0.04\times(8+0.10\times90)=1.68\ \mathrm{ns}.
\]
真正到达主存的比例为 $0.04\times0.10=0.4\%$。
\answercheck{必须掌握的要点}{下一级局部 miss 率只以到达该级的请求为分母}
\end{answersubproblem}
\begin{answersubproblem}{（b）冲突 trace}
2 KiB/64 B 得 32 个集合；三个地址分别相差 2 KiB，line 编号相差 32，取模后都进入集合 0。A、B、C、A 全部 miss，每次都替换该集合唯一的一路。
\answercheck{必须掌握的要点}{工作集小于容量也可能因映射集中而持续 miss}
\end{answersubproblem}
\end{chapteranswer}

\begin{chapteranswer}{综合练习：暂态闭合至目录事务}
\begin{answersubproblem}{（a）暂态闭合}
不能。收到数据只满足数据条件，尚未取得唯一写权限；必须等待两个确认归零后才能进入稳定 M 并让写生效。暂态要保存目标状态、待确认计数、数据返回状态、请求者和错误标记。
\answercheck{必须掌握的要点}{权限完成不能由数据到达替代}
\end{answersubproblem}
\begin{answersubproblem}{（b）目录事务}
C1 向 home 发共享读；home 查到 C0 是 owner，转发请求；C0 把最新数据发送给 C1，并向 home 提供状态更新，从 M 降到 S；home 把 C0、C1 记为共享者，C1 收到数据后进入 S。内存副本在 C0 为 M 时可能陈旧，不能直接使用。
\answercheck{必须掌握的要点}{请求要到达持有最新 dirty 数据的 owner}
\end{answersubproblem}
\end{chapteranswer}

\begin{chapteranswer}{综合练习：伪共享预算至自修改代码}
\begin{answersubproblem}{（a）伪共享预算}
所有权每 100 ns 最多完成一次迁移，合计更新率上限约为 $1/100\ \mathrm{ns}=1000$ 万次/秒，低于两个本地更新器的需求。按 line 分离后，各核心可长期持有自己的 M line，更新主要受本地 L1 和执行端限制，只在汇总时通信。
\answercheck{必须掌握的要点}{瓶颈从 ALU/本地 Cache 转为 line 所有权往返}
\end{answersubproblem}
\begin{answersubproblem}{（b）自修改代码}
先使数据写到达指令侧可见的一致性点，再失效可能保存旧字节的 I-Cache 项，最后执行指令同步操作，丢弃流水线、预取队列或译码缓存中的旧指令。三个边界缺一不可。
\answercheck{必须掌握的要点}{D-Cache、I-Cache 和前端是三个不同状态持有者}
\end{answersubproblem}
\end{chapteranswer}

\begin{chapteranswer}{综合练习：SRAM 读写周期至 cache 行}
\begin{answersubproblem}{（a）SRAM 读写周期}
读时地址和片选稳定，OE 有效、WE 无效，由 SRAM 驱动数据；写时 OE 无效，由发起者驱动数据，WE 在地址和数据稳定后有效。
\answercheck{必须保留的要点}{输出使能互斥，写窗口满足建立/保持时间。}
\end{answersubproblem}
\begin{answersubproblem}{（b）cache 行}
先算 offset=log2（行大小）、index=log2（行数）、tag=剩余位；命中要求槽 valid 且 tag 相等；miss 时取整行填入并更新 tag/valid，dirty 位决定替换时是否要先写回下一级。
\answercheck{必须保留的要点}{三个字段的分工和 valid/dirty 语义不能丢。}
\end{answersubproblem}
\end{chapteranswer}

\begin{chapteranswer}{综合练习：组相联查找至替换策略}
\begin{answersubproblem}{（a）组相联查找}
index 选中整个组，组内两路 tag 并行与地址 tag 比较，任一路相等且 V=1 即 hit；两个热地址各占一路后不再互相驱逐；miss 时按替换位选一路驱逐。
\answercheck{必须保留的要点}{并行比较、mux 和替换位缺一不可。}
\end{answersubproblem}
\begin{answersubproblem}{（b）替换策略}
2 路 LRU：A miss 得 \{A,--\}；B miss 得 \{A,B\}；A hit（B 成为最久未用）；C miss 驱逐 B 得 \{A,C\}；A hit；B miss 驱逐 C 得 \{A,B\}，共 3 次 miss。直接映射下同 index 只有一个槽，每次访问都驱逐上一行，6 次全 miss。
\answercheck{必须保留的要点}{LRU 驱逐“最久未用”，不是“最早进入”。}
\end{answersubproblem}
\end{chapteranswer}

\begin{chapteranswer}{预取演算}
无预取：每行 miss 30 拍加处理 8 拍，$16\times38=608$ 拍。next-line 预取：第 1 行第 30 拍就绪，处理第 $i$ 行的 8 拍只能隐藏预取的一部分，第 $i$ 行数据在第 $30i$ 拍就绪，稳态每行 30 拍，总 $30\times16+8=488$ 拍。处理节拍 $T\geq30$ 时预取被完全隐藏，总拍数降为 $30+16T$。
\answercheck{必须保留的要点}{预取隐藏延迟但不消除首次 miss；处理太慢时仍受 miss 间隔限制。}
\end{chapteranswer}

\begin{chapteranswer}{MESI 走查}
步骤 1：A 由 I 变 S，B 保持 I（读入共享副本）；步骤 2：B 由 I 变 S，A 保持 S；步骤 3：A 广播失效，B 由 S 变 I，A 由 S 变 M，内存过期；步骤 4：B 读 miss，A 写回或转发后由 M 降为 S，B 得 S，内存恢复最新。
\answercheck{必须保留的要点}{每步最新值唯一：S 期间在内存，M 期间在 A 的 cache。}
\end{chapteranswer}
